% =====================================================================
%  SECCION 5 -- PRIORIZACION Y TRAZABILIDAD EXTENDIDA
%  Se incluye desde ERS_SRS_2A_v1.0.tex mediante % =====================================================================
%  SECCION 5 -- PRIORIZACION Y TRAZABILIDAD EXTENDIDA
%  Se incluye desde ERS_SRS_2A_v1.0.tex mediante % =====================================================================
%  SECCION 5 -- PRIORIZACION Y TRAZABILIDAD EXTENDIDA
%  Se incluye desde ERS_SRS_2A_v1.0.tex mediante % =====================================================================
%  SECCION 5 -- PRIORIZACION Y TRAZABILIDAD EXTENDIDA
%  Se incluye desde ERS_SRS_2A_v1.0.tex mediante \input{seccion5_priorizacion}
% =====================================================================

\newpage
\section{Priorización y trazabilidad extendida}

\subsection{Clasificación de requisitos}

El documento especifica 75 elementos de requisito, distribuidos como sigue.

\begin{longtable}{|L{4.2cm}|C{2.0cm}|L{8.4cm}|}
\hline
\rowcolor{grisclaro}\textbf{Tipo} & \textbf{Cantidad} & \textbf{Identificadores} \\
\hline
\endhead
Requisito funcional (RF) & 39 & \id{RF-01} a \id{RF-39} \\ \hline
Requisito no funcional (RNF) & 18 & \id{RNF-01} a \id{RNF-18} \\ \hline
Restricción de diseño (RD) & 10 & \id{RD-01} a \id{RD-10} \\ \hline
Requisito legal (RL) & 8 & \id{RL-01} a \id{RL-08} \\ \hline
\rowcolor{grisclaro}\textbf{Total} & \textbf{75} & \\ \hline
\caption{Clasificación de los requisitos por tipo}
\end{longtable}

\noindent
La distribución MoSCoW de los requisitos funcionales es de 21 \emph{Must}, 16 \emph{Should} y 2
\emph{Could}. Los cuatro requisitos incorporados a partir de la hoja de liquidación semanal
(\id{RF-36} a \id{RF-39}) se clasifican como \emph{Must} los dos primeros ---el catálogo de tarifas
es prerrequisito del cálculo de remuneración y la liquidación es el cierre del ciclo semanal--- y
como \emph{Should} los dos restantes. No hay requisitos funcionales clasificados como \emph{Won't}: las exclusiones de
alcance se enuncian a nivel de producto en la §1.2.4 y se argumentan en la §5.3.4.

% ---------------------------------------------------------------------
\subsection{Historias de usuario de los requisitos \emph{Must}}

Conforme a la §3.3 de la guía de la Entrega 3, cada requisito funcional de prioridad \emph{Must}
cuenta con una historia de usuario en formato Connextra, con los criterios INVEST verificados y un
escenario de aceptación redactado en Gherkin. Se especifican 19 historias (\id{HU-01} a
\id{HU-19}), con sus criterios de aceptación (\id{CA-01} a \id{CA-19}).

\begin{longtable}{|C{1.3cm}|C{1.3cm}|L{12.4cm}|}
\hline
\rowcolor{grisclaro}\textbf{HU} & \textbf{RF} & \textbf{Historia de usuario (Connextra) y escenario de aceptación (Gherkin)} \\
\hline
\endhead

\id{HU-01} & \id{RF-01} & \textbf{Como} persona trabajadora del cultivo, \textbf{quiero} acceder al sistema con mis credenciales, \textbf{para} usar únicamente las funciones que corresponden a mi labor. \newline \emph{\id{CA-01}:} \textbf{Dado} que existe una cuenta con rol ``trabajador agrícola'', \textbf{cuando} inicia sesión, \textbf{entonces} el menú no presenta ninguna opción de administración \textbf{y} un intento de acceso directo a esa ruta es rechazado. \\ \hline

\id{HU-02} & \id{RF-02} & \textbf{Como} administrador general, \textbf{quiero} registrar los lotes con su variedad y superficie, \textbf{para} poder asociarles labores, diagnósticos y cosechas. \newline \emph{\id{CA-02}:} \textbf{Dado} un lote recién registrado con variedad asignada, \textbf{cuando} se abre el formulario de registro de labor, \textbf{entonces} el lote aparece como opción seleccionable. \\ \hline

\id{HU-03} & \id{RF-03} & \textbf{Como} administrador general, \textbf{quiero} organizar al personal por equipos funcionales, \textbf{para} asignar labores y evaluar el desempeño por cuadrilla. \newline \emph{\id{CA-03}:} \textbf{Dado} un trabajador asignado al equipo de polinización, \textbf{cuando} se consulta el equipo, \textbf{entonces} la persona figura en él \textbf{y} es asignable a labores de ese tipo. \\ \hline

\id{HU-04} & \id{RF-04} & \textbf{Como} trabajador agrícola, \textbf{quiero} registrar la labor que ejecuté en el lote, \textbf{para} que quede constancia sin depender de una libreta de papel. \newline \emph{\id{CA-04}:} \textbf{Dado} que ejecuté chapia en el lote 3, \textbf{cuando} registro la labor con su cantidad, \textbf{entonces} queda vinculada al lote, la fecha y mi identificación \textbf{y} aparece en el historial del lote. \\ \hline

\id{HU-05} & \id{RF-05} & \textbf{Como} técnico fitosanitario, \textbf{quiero} registrar el monitoreo sanitario del lote, \textbf{para} dejar trazabilidad de las revisiones y de las incidencias detectadas. \newline \emph{\id{CA-05}:} \textbf{Dado} un monitoreo con incidencia marcada, \textbf{cuando} se filtran los lotes con incidencia abierta, \textbf{entonces} ese lote aparece en el resultado. \\ \hline

\id{HU-06} & \id{RF-07} & \textbf{Como} asistente de administración, \textbf{quiero} fotografiar una palma afectada y obtener el diagnóstico, \textbf{para} identificar la plaga sin depender de conocimiento especializado. \newline \emph{\id{CA-06}:} \textbf{Dado} que capturo la imagen de una palma con afectación conocida, \textbf{cuando} solicito el análisis, \textbf{entonces} el sistema devuelve la clase, el nivel de confianza y una explicación de no más de 60 palabras \textbf{y} registra el diagnóstico con lote y fecha. \\ \hline

\id{HU-07} & \id{RF-08} & \textbf{Como} asesor técnico, \textbf{quiero} fotografiar una hoja y conocer la deficiencia nutricional, \textbf{para} decidir qué aplicar sin esperar el análisis de laboratorio. \newline \emph{\id{CA-07}:} \textbf{Dado} el lote con variedad asignada, \textbf{cuando} envío la imagen de una hoja con deficiencia conocida, \textbf{entonces} el sistema indica el elemento deficiente, propone una corrección \textbf{y} explica qué signo foliar la sustenta. \\ \hline

\id{HU-08} & \id{RF-10} & \textbf{Como} asesor técnico, \textbf{quiero} que el sistema conozca la variedad de cada lote, \textbf{para} que los umbrales de diagnóstico correspondan al material sembrado. \newline \emph{\id{CA-08}:} \textbf{Dado} un lote de variedad \emph{guineensis}, \textbf{cuando} se clasifica un racimo, \textbf{entonces} el umbral aplicado es de 3 frutos desprendidos; \textbf{y dado} un lote híbrido, \textbf{entonces} el umbral aplicado es de 5. \\ \hline

\id{HU-09} & \id{RF-12} & \textbf{Como} administrador general, \textbf{quiero} recibir alertas ante una afectación detectada, \textbf{para} intervenir antes de que se propague. \newline \emph{\id{CA-09}:} \textbf{Dado} un diagnóstico positivo registrado, \textbf{cuando} se completa el análisis, \textbf{entonces} se crea una alerta con lote, tipo, criticidad y estado abierto \textbf{y} resulta visible para el rol responsable. \\ \hline

\id{HU-10} & \id{RF-13} & \textbf{Como} jefe de polinización, \textbf{quiero} registrar el estado de la antesis y la verificación visual, \textbf{para} controlar la calidad de la labor que determina la producción. \newline \emph{\id{CA-10}:} \textbf{Dado} un lote en etapa de polinización, \textbf{cuando} se registra la jornada con estado de antesis, \textbf{entonces} el registro queda asociado al lote, la fecha y la persona \textbf{y} es consultable por fecha. \\ \hline

\id{HU-11} & \id{RF-14} & \textbf{Como} polinizador, \textbf{quiero} que el conteo de flores se calcule solo, \textbf{para} no perder rendimiento anotando planta por planta. \newline \emph{\id{CA-11}:} \textbf{Dado} que trabajé la jornada con rastreo activo, \textbf{cuando} cierro la jornada, \textbf{entonces} el sistema calcula el total de flores a partir de las marcaciones \textbf{y} no me solicita ningún registro manual por planta. \\ \hline

\id{HU-12} & \id{RF-18} & \textbf{Como} administrador general, \textbf{quiero} estimar la producción del período, \textbf{para} planificar la cosecha y la mano de obra necesaria. \newline \emph{\id{CA-12}:} \textbf{Dado} un censo de flores registrado y la variedad del lote, \textbf{cuando} solicito la estimación, \textbf{entonces} el sistema devuelve el tonelaje estimado del período \textbf{y} el supuesto de rendimiento empleado. \\ \hline

\id{HU-13} & \id{RF-19} & \textbf{Como} administrador general, \textbf{quiero} generar el consolidado semanal, \textbf{para} tomar decisiones sin recopilar reportes dispersos de mensajería. \newline \emph{\id{CA-13}:} \textbf{Dado} que existen registros en la semana, \textbf{cuando} genero el reporte, \textbf{entonces} obtengo labores, avance, producción y alertas del período \textbf{y} puedo exportarlo a PDF o a hoja de cálculo. \\ \hline

\id{HU-14} & \id{RF-21} & \textbf{Como} cosechador, \textbf{quiero} saber si un racimo está en punto de corte, \textbf{para} no cortar fruta verde que será penalizada. \newline \emph{\id{CA-14}:} \textbf{Dado} un lote híbrido, \textbf{cuando} fotografío la parte basal de un racimo con 4 frutos desprendidos, \textbf{entonces} el sistema lo clasifica como verde \textbf{y} explica que el umbral de la variedad es de 5. \\ \hline

\id{HU-15} & \id{RF-22} & \textbf{Como} administrador general, \textbf{quiero} conocer el porcentaje estimado de fruta verde antes de despachar, \textbf{para} evitar la penalización del lote en la planta extractora. \newline \emph{\id{CA-15}:} \textbf{Dado} un umbral de penalización del 3\%, \textbf{cuando} el porcentaje estimado del lote alcanza ese valor, \textbf{entonces} el sistema emite alerta preventiva \textbf{y} desaconseja el despacho. \\ \hline

\id{HU-16} & \id{RF-26} & \textbf{Como} administrador general, \textbf{quiero} registrar el presupuesto semanal de labores, \textbf{para} contrastar lo planificado con lo efectivamente ejecutado. \newline \emph{\id{CA-16}:} \textbf{Dado} un plan con meta de 1.123 coronas en el lote 3, \textbf{cuando} cierro la semana con 800 cumplidas, \textbf{entonces} el sistema reporta un cumplimiento del 71,2\% \textbf{y} un remanente de 323 coronas. \\ \hline

\id{HU-17} & \id{RF-28} & \textbf{Como} trabajador agrícola, \textbf{quiero} registrar mi avance en la unidad propia de cada labor, \textbf{para} que mi trabajo se contabilice como corresponde. \newline \emph{\id{CA-17}:} \textbf{Dado} que ejecuté 120 coronas, \textbf{cuando} registro el avance, \textbf{entonces} la cantidad se acumula a mi total semanal en la unidad ``coronas'' \textbf{y} el sistema rechaza una unidad incompatible con el tipo de labor. \\ \hline

\id{HU-18} & \id{RF-30} & \textbf{Como} asistente de administración, \textbf{quiero} reportar una afectación con foto y ubicación, \textbf{para} poder volver al punto exacto en mi siguiente visita. \newline \emph{\id{CA-18}:} \textbf{Dado} que reporto una incidencia con GPS activo, \textbf{cuando} la consulto días después, \textbf{entonces} recupero su fotografía y sus coordenadas \textbf{y} la incidencia figura en mi lista de pendientes de verificación. \\ \hline

\id{HU-19} & \id{RF-35} & \textbf{Como} jefe de campo, \textbf{quiero} registrar el avance de quien no tiene teléfono, \textbf{para} que ninguna persona quede sin registrar su trabajo. \newline \emph{\id{CA-19}:} \textbf{Dado} un trabajador sin dispositivo propio, \textbf{cuando} registro su avance de forma delegada, \textbf{entonces} el avance se atribuye a esa persona en su consolidado \textbf{y} la bitácora conserva mi identificación como registrador. \\ \hline
\caption{Historias de usuario y criterios de aceptación de los requisitos \emph{Must}}
\end{longtable}

\subsubsection{Verificación de los criterios INVEST}

Las diecinueve historias fueron evaluadas contra los seis criterios INVEST. La tabla siguiente
reporta el resultado y documenta las tres historias que requirieron ajuste.

\begin{longtable}{|L{2.4cm}|C{1.6cm}|L{11.0cm}|}
\hline
\rowcolor{grisclaro}\textbf{Criterio} & \textbf{Cumple} & \textbf{Observación} \\
\hline
\endhead
Independiente & 16 / 19 & \id{HU-11} depende de \id{HU-10}, \id{HU-15} de \id{HU-14} y \id{HU-17} de \id{HU-04}. Se documenta la dependencia en la matriz en lugar de forzar la independencia, dado que responde a la secuencia real del proceso. \\ \hline
Negociable & 19 / 19 & Ninguna historia prescribe solución técnica; los umbrales son parametrizables (\id{RD-08}). \\ \hline
Valiosa & 19 / 19 & Cada historia declara el beneficio esperado y se traza a evidencia de campo. \\ \hline
Estimable & 19 / 19 & Todas admiten estimación en la escala Fibonacci empleada en el cálculo WSJF de la §5.3.3. \\ \hline
Pequeña & 17 / 19 & \id{HU-06} y \id{HU-07} superan el tamaño recomendado por incluir el componente de IA; se sugiere descomponerlas en captura, inferencia y explicación durante la planificación de la implementación. \\ \hline
Verificable & 19 / 19 & Todas cuentan con escenario Gherkin con criterio observable. \\ \hline
\caption{Verificación de los criterios INVEST sobre las historias de usuario}
\end{longtable}

% ---------------------------------------------------------------------
\subsection{Priorización combinada}

La priorización emplea tres instrumentos complementarios: MoSCoW para el compromiso de alcance, el
modelo de Kano para la relación con la satisfacción del cliente, y WSJF para el orden de
implementación.

\subsubsection{Justificación de la clasificación MoSCoW}

\textbf{\emph{Must} (19 RF).} Corresponden a tres condiciones. Primero, la \textbf{base operativa
indispensable}: sin autenticación por rol (\id{RF-01}), estructura de lotes (\id{RF-02}), personal
(\id{RF-03}) y registro de labor y avance (\id{RF-04}, \id{RF-28}, \id{RF-35}) el sistema no
sustituye al proceso manual que pretende reemplazar. Segundo, el \textbf{diferencial de valor}: el
análisis por imagen (\id{RF-07}, \id{RF-08}, \id{RF-21}) y las alertas (\id{RF-12}) son las
funciones mejor valoradas por el personal en el cuestionario ($4{,}31$ y $4{,}27$ sobre 5,
\id{EV-12}) y las señaladas de forma independiente por la administración (\id{EV-01}) y la asesoría
técnica (\id{EV-02}). Tercero, la \textbf{consecuencia económica directa}: la polinización determina
la producción de la finca (\id{EV-03}), su conteo condiciona la evaluación del trabajo
(\id{RF-13}, \id{RF-14}), y el porcentaje de fruta verde determina la penalización del lote en la
planta extractora (\id{RF-22}, \id{EV-04}).

\textbf{\emph{Should} (14 RF).} Aportan valor sustantivo pero el sistema entrega su propuesta de
valor sin ellos en la primera versión: seguimiento de etapas, ficha de plaga, análisis de
suelo, recorridos, evaluación de desempeño, clima, historial, registro del ticket de la extractora,
trazabilidad de la penalización, ventana corte-entrega, arrastre semanal, consulta de remuneración,
lista de pendientes y exportación hacia la extractora.

\textbf{\emph{Could} (2 RF).} \id{RF-33} (guía de remisión) y \id{RF-34} (solicitud de insumos) son
mejoras deseables de bajo acoplamiento con el núcleo.

\textbf{\emph{Won't} en esta versión.} Se excluyen explícitamente y con argumento: la
\textbf{gestión de nómina y pago}, porque el sistema calcula el avance pero la liquidación
involucra información laboral y tributaria cuyo tratamiento excede el consentimiento obtenido
(\id{RD-04}, \id{RL-02}); la \textbf{facturación y contabilidad}, por no corresponder al manejo del
cultivo; la \textbf{comercialización de la producción}, que se negocia fuera del sistema; la
\textbf{integración directa con el sistema de la planta extractora}, porque pertenece a una
organización externa con sistema propio y su interfaz no está disponible (\id{EV-04}); y la
\textbf{operación multi-finca}, que el cliente administra pero que no forma parte del alcance
acordado para la v1 (§1.2.2).

\subsubsection{Modelo de Kano}

La clasificación se sustenta en la distancia entre la valoración declarada en el cuestionario y la
expectativa implícita del proceso actual.

\begin{longtable}{|L{2.8cm}|L{5.4cm}|L{7.0cm}|}
\hline
\rowcolor{grisclaro}\textbf{Categoría} & \textbf{Requisitos} & \textbf{Fundamento} \\
\hline
\endhead
Básica & \id{RF-01}, \id{RF-02}, \id{RF-03}, \id{RF-04} & Su ausencia genera rechazo, su presencia no genera entusiasmo. El registro de labor ya existe en papel: el sistema debe igualarlo, no impresiona por hacerlo. \\ \hline
De desempeño & \id{RF-19}, \id{RF-26}, \id{RF-28}, \id{RF-29}, \id{RF-13} & La satisfacción crece de forma proporcional a la calidad de ejecución. Valoración intermedia en el cuestionario ($3{,}95$ para el registro de labores). \\ \hline
\rowcolor{grisclaro}De deleite & \id{RF-07}, \id{RF-08}, \id{RF-12}, \id{RF-21} & Generan entusiasmo desproporcionado respecto de la expectativa. Son las funciones mejor valoradas ($4{,}27$ y $4{,}31$) y las únicas mencionadas de forma espontánea por los cinco perfiles entrevistados. \\ \hline
Indiferente & \id{RF-33}, \id{RF-34} & No inciden en la satisfacción percibida por la persona usuaria final. \\ \hline
De rechazo potencial & \id{RF-15} & El rastreo de recorridos puede generar insatisfacción: el 25{,}8\% del personal manifiesta reservas (\id{EV-12}). Se mitiga con consentimiento revocable (\id{RL-03}). \\ \hline
\caption{Clasificación de los requisitos según el modelo de Kano}
\end{longtable}

\noindent
La clasificación de \id{RF-15} como \emph{requisito de rechazo potencial} es deliberada y
constituye un hallazgo. La funcionalidad fue solicitada por la administración (\id{EV-01},
\id{EV-02}) y por la jefatura de polinización (\id{EV-03}), pero una de cada cuatro personas
destinatarias expresa reservas sobre ella. Kano permite representar esta tensión: un requisito puede
ser valioso para quien lo solicita y adverso para quien lo padece. La mitigación no es técnica sino
de diseño de la relación: consentimiento específico, revocable y sin consecuencia laboral.

\subsubsection{Valor de negocio mediante WSJF}

Se aplica el método del Trabajo Ponderado Más Corto Primero:
\[
\mathrm{WSJF} = \frac{\text{Valor de negocio} + \text{Criticidad temporal} + \text{Reducción de riesgo}}{\text{Tamaño del trabajo}}
\]
Los cuatro componentes se puntúan en escala Fibonacci (1, 2, 3, 5, 8, 13). La tabla presenta los
diez requisitos evaluados de mayor valor de negocio, ordenados de forma descendente por WSJF.

\begin{longtable}{|C{1.4cm}|C{1.3cm}|C{1.3cm}|C{1.3cm}|C{1.3cm}|C{1.4cm}|L{5.4cm}|}
\hline
\rowcolor{grisclaro}\textbf{RF} & \textbf{Valor} & \textbf{Crit. temp.} & \textbf{Red. riesgo} & \textbf{Tamaño} & \textbf{WSJF} & \textbf{Justificación} \\
\hline
\endhead
\id{RF-35} & 5 & 8 & 8 & 3 & \textbf{7,00} & Bajo costo y alto impacto: sin él, el 11,3\% del personal queda excluido del sistema (\id{EV-12}). \\ \hline
\id{RF-21} & 13 & 13 & 8 & 8 & \textbf{4,25} & Evita la penalización económica del lote; la temporada de baja producción (jul.--oct.) es cuando más fruta verde se corta (\id{EV-04}). \\ \hline
\id{RF-22} & 13 & 13 & 8 & 8 & \textbf{4,25} & Actúa antes del despacho, momento en que la pérdida aún es evitable. \\ \hline
\id{RF-04} & 8 & 8 & 5 & 5 & \textbf{4,20} & Base de todo el registro operativo; sin él no hay datos para ninguna otra función. \\ \hline
\id{RF-12} & 8 & 8 & 5 & 5 & \textbf{4,20} & Función mejor valorada del cuestionario (4,31/5). \\ \hline
\id{RF-01} & 5 & 8 & 8 & 5 & \textbf{4,20} & Precondición técnica y legal de todo lo demás (\id{RL-01}). \\ \hline
\id{RF-07} & 13 & 8 & 8 & 8 & \textbf{3,62} & Mayor dificultad declarada por el personal (40,3\%, \id{EV-12}). \\ \hline
\id{RF-28} & 8 & 5 & 5 & 5 & \textbf{3,60} & Habilita el cálculo de cumplimiento y de remuneración por avance. \\ \hline
\id{RF-14} & 13 & 8 & 5 & 8 & \textbf{3,25} & La polinización determina la producción; el conteo manual afecta el rendimiento (\id{EV-03}). \\ \hline
\id{RF-08} & 8 & 5 & 5 & 8 & \textbf{2,25} & Alto valor técnico, pero el diagnóstico nutricional admite mayor demora que el sanitario. \\ \hline

\caption{Valor de negocio mediante WSJF (diez requisitos de mayor prioridad)}
\end{longtable}

\noindent
\textbf{Lectura del resultado.} \id{RF-35} y \id{RF-10} encabezan el orden de implementación con un
WSJF de 7,00 cada uno,
pese a no ser el requisito de mayor valor de negocio. La razón es su tamaño reducido frente a un
impacto alto: se trata de una variante del formulario de registro que evita excluir a una de cada
nueve personas destinatarias. \id{RF-10} alcanza la misma posición por una razón distinta: es
prerrequisito de \id{RF-21}, dado que el umbral de desprendimiento se obtiene de la variedad del
lote. El método revela así una dependencia técnica que ni MoSCoW ni Kano hacen visible.

Ambos casos ilustran la utilidad del método frente a una priorización basada únicamente en valor
percibido, que los habría situado muy por debajo. La lectura inversa también es informativa:
\id{RF-07} y \id{RF-08}, las dos funciones mejor valoradas por las personas usuarias, obtienen WSJF
de 3,62 y 2,25 por su elevado tamaño de trabajo. Esto no altera su condición de \emph{Must}, pero
indica que el desarrollo debe iniciarse por la base operativa mientras los componentes de
inteligencia artificial se construyen en paralelo. El archivo
\id{04\_Trazabilidad/}\allowbreak\id{priorizacion\_moscow\_kano.csv} contiene la tabla completa para los 35
requisitos funcionales.

% ---------------------------------------------------------------------
\subsection{Matriz de trazabilidad extendida}

La matriz enlaza, para cada elemento, la cadena completa
Ley $\rightarrow$ Objetivo $\rightarrow$ Interesado $\rightarrow$ Evidencia $\rightarrow$
Requisito $\rightarrow$ Caso de uso $\rightarrow$ Historia de usuario $\rightarrow$
Criterio de aceptación $\rightarrow$ Componente $\rightarrow$ Mockup.

La versión completa, con las trece columnas exigidas por la §8.8 de la guía, reside en
\id{04\_Trazabilidad/}\allowbreak\id{matriz\_trazabilidad.csv} y contiene \textbf{52 filas}. La tabla siguiente
reproduce una vista reducida con las columnas de mayor densidad informativa; se conservan las 48
filas para permitir la verificación de cobertura. Las cuatro últimas corresponden a los requisitos derivados de la hoja de liquidación semanal (\id{EV-13}).

\begin{longtable}{|C{0.9cm}|C{1.5cm}|C{1.6cm}|C{1.5cm}|C{1.3cm}|C{1.3cm}|C{1.3cm}|C{1.5cm}|}
\hline
\rowcolor{grisclaro}\textbf{\#} & \textbf{Interesado} & \textbf{ID-EV} & \textbf{ID-RF} & \textbf{Tipo} & \textbf{ID-CU} & \textbf{ID-HU} & \textbf{ID-Mockup} \\
\hline
\endhead
1 & Administrador & EV-01, EV-09 & RF-01 & RF & CU-11 & HU-01 & MU-01 \\ \hline
2 & Administrador & EV-01, EV-02 & RF-02 & RF & CU-12 & HU-02 & MU-05 \\ \hline
3 & Administrador & EV-01, EV-11 & RF-03 & RF & CU-13 & HU-03 & --- \\ \hline
4 & Trabajador & EV-05, EV-06 & RF-04 & RF & CU-04 & HU-04 & MU-04 \\ \hline
5 & Trabajador & EV-11 & RF-04 & RF & CU-04 & HU-04 & MU-04 \\ \hline
6 & Técnico fitosan. & EV-01, EV-02 & RF-05 & RF & CU-15 & HU-05 & MU-05 \\ \hline
7 & Administrador & EV-01 & RF-06 & RF & CU-12 & --- & MU-05 \\ \hline
8 & Asistente admin. & EV-01, EV-02 & RF-07 & RF & CU-01 & HU-06 & MU-03 \\ \hline
9 & Asistente admin. & EV-07 & RF-07 & RF & CU-01 & HU-06 & MU-03 \\ \hline
10 & Asesor técnico & EV-02 & RF-08 & RF & CU-02 & HU-07 & MU-03 \\ \hline
11 & Asesor técnico & EV-04 & RF-08 & RF & CU-02 & HU-07 & MU-03 \\ \hline
12 & Técnico fitosan. & EV-01, EV-02 & RF-09 & RF & CU-01 & --- & MU-03 \\ \hline
13 & Asesor técnico & EV-04 & RF-10 & RF & CU-12 & HU-08 & MU-05 \\ \hline
14 & Asesor técnico & EV-02 & RF-11 & RF & CU-02 & --- & MU-05 \\ \hline
15 & Administrador & EV-01, EV-02 & RF-12 & RF & CU-16 & HU-09 & MU-08 \\ \hline
16 & Trabajador & EV-12 & RF-12 & RF & CU-16 & HU-09 & MU-08 \\ \hline
17 & Jefe polinización & EV-03 & RF-13 & RF & CU-03 & HU-10 & MU-04 \\ \hline
18 & Jefe polinización & EV-03 & RF-14 & RF & CU-03 & HU-11 & MU-06 \\ \hline
19 & Administrador & EV-01, EV-03 & RF-15 & RF & CU-03 & --- & MU-06 \\ \hline
20 & Administrador & EV-01, EV-02 & RF-16 & RF & CU-05 & --- & MU-07 \\ \hline
21 & Administrador & EV-01, EV-02 & RF-17 & RF & CU-05 & --- & MU-05 \\ \hline
22 & Administrador & EV-01, EV-02 & RF-18 & RF & CU-17 & HU-12 & MU-07 \\ \hline
23 & Administrador & EV-01, EV-11 & RF-19 & RF & CU-05 & HU-13 & MU-07 \\ \hline
24 & Administrador & EV-01 & RF-20 & RF & CU-05 & --- & MU-05 \\ \hline
25 & Téc. extractora & EV-04 & RF-21 & RF & CU-06 & HU-14 & MU-03 \\ \hline
26 & Téc. extractora & EV-08 & RF-21 & RF & CU-06 & HU-14 & MU-03 \\ \hline
27 & Administrador & EV-04, EV-08 & RF-22 & RF & CU-06 & HU-15 & MU-08 \\ \hline
28 & Administrador & EV-04 & RF-23 & RF & CU-10 & --- & MU-07 \\ \hline
29 & Téc. extractora & EV-04, EV-08 & RF-24 & RF & CU-10 & --- & MU-07 \\ \hline
30 & Téc. extractora & EV-04 & RF-25 & RF & CU-10 & --- & MU-07 \\ \hline
31 & Administrador & EV-11 & RF-26 & RF & CU-07 & HU-16 & MU-02 \\ \hline
32 & Administrador & EV-11 & RF-27 & RF & CU-07 & --- & MU-02 \\ \hline
33 & Trabajador & EV-05, EV-11 & RF-28 & RF & CU-04 & HU-17 & MU-04 \\ \hline
34 & Trabajador & EV-06, EV-07 & RF-29 & RF & CU-08 & --- & MU-02 \\ \hline
35 & Asistente admin. & EV-07 & RF-30 & RF & CU-09 & HU-18 & MU-06 \\ \hline
36 & Asistente admin. & EV-07 & RF-31 & RF & CU-09 & --- & MU-08 \\ \hline
37 & Téc. extractora & EV-04 & RF-32 & RF & CU-17 & --- & MU-07 \\ \hline
38 & Téc. extractora & EV-04 & RF-33 & RF & CU-10 & --- & MU-07 \\ \hline
39 & Administrador & EV-11 & RF-34 & RF & CU-07 & --- & MU-02 \\ \hline
40 & Jefe de campo & EV-12, EV-05 & RF-35 & RF & CU-04 & HU-19 & MU-04 \\ \hline
41 & Téc. extractora & EV-04 & RNF-01 & RNF & CU-06 & HU-14 & MU-03 \\ \hline
42 & Trabajador & EV-07, EV-12 & RNF-11 & RNF & CU-04 & HU-04 & MU-04 \\ \hline
43 & Jefe polinización & EV-03 & RNF-05 & RNF & CU-03 & HU-11 & MU-06 \\ \hline
44 & Trabajador & EV-12 & RNF-08 & RNF & CU-04 & HU-04 & MU-04 \\ \hline
45 & Trabajador & EV-12 & RNF-14 & RNF & CU-03 & HU-11 & MU-06 \\ \hline
46 & Asistente admin. & EV-01, EV-07 & RNF-16 & RNF & CU-18 & HU-06 & MU-03 \\ \hline
47 & Trabajador & EV-12 & RD-10 & RD & CU-04 & HU-19 & MU-04 \\ \hline
48 & Téc. extractora & EV-08 & RD-07 & RD & CU-06 & HU-14 & MU-03 \\ \hline
49 & Administrador & EV-13 & RF-36 & RF & CU-08 & --- & MU-02 \\ \hline
50 & Administrador & EV-13, EV-11 & RF-37 & RF & CU-07 & --- & MU-02 \\ \hline
51 & Administrador & EV-13 & RF-38 & RF & CU-07 & --- & MU-02 \\ \hline
52 & Administrador & EV-13 & RF-39 & RF & CU-07 & --- & MU-02 \\ \hline
\caption{Matriz de trazabilidad extendida (vista reducida, 52 filas)}
\end{longtable}

\subsubsection{Diagnóstico de cobertura}

\begin{longtable}{|L{5.6cm}|C{2.4cm}|L{6.8cm}|}
\hline
\rowcolor{grisclaro}\textbf{Indicador} & \textbf{Resultado} & \textbf{Interpretación} \\
\hline
\endhead
Requisitos funcionales trazados a evidencia & 39 / 39 (100\%) & Ningún requisito carece de fuente verificable. \\ \hline
Requisitos \emph{Must} con caso de uso & 19 / 19 (100\%) & Se corrige la deficiencia de la Entrega 2, donde \id{RF-17} y \id{RF-20} carecían de caso de uso. \\ \hline
Requisitos \emph{Must} con historia y criterio & 19 / 19 (100\%) & \\ \hline
Evidencias que originan al menos un requisito & 12 / 13 (92,3\%) & \id{EV-09} (observación de campo) aporta contexto y respalda requisitos junto con otras fuentes, sin originar ninguno de forma exclusiva. \\ \hline
Requisitos sustentados en una sola evidencia & 9 / 39 (23,1\%) & Véase el análisis de elementos huérfanos. \\ \hline
Requisitos con mockup asociado & 36 / 39 (92,3\%) & \id{RF-03}, \id{RF-27} y \id{RF-34} operan sobre pantallas no prototipadas en esta entrega. \\ \hline
\caption{Diagnóstico de cobertura de la matriz de trazabilidad}
\end{longtable}

\subsubsection{Elementos huérfanos y vacíos de cobertura}

Se declaran de forma explícita las debilidades detectadas, conforme al criterio de completitud de la
guía.

\begin{enumerate}[nosep]
  \item \textbf{Requisitos con fuente única.} \id{RF-11}, \id{RF-23}, \id{RF-25}, \id{RF-32},
        \id{RF-33} y \id{RF-34} se sustentan en una sola evidencia. Los cuatro derivados de
        \id{EV-04} conciernen a la relación con la planta extractora y no pudieron triangularse con
        una segunda fuente independiente de esa organización, dado que \id{EV-08} corresponde a la
        misma sesión de campo. Se declara como limitación.
  \item \textbf{Requisitos sin prototipo de interfaz.} \id{RF-03}, \id{RF-27} y \id{RF-34}. Los tres
        son de prioridad no \emph{Must}, por lo que la ausencia no incumple la exigencia de
        prototipar las pantallas \emph{Must}.
  \item \textbf{Ausencia de \emph{Won't} funcionales.} Las exclusiones se enuncian a nivel de
        producto y no como requisitos individuales descartados. Se argumentan en la §5.3.1.
  \item \textbf{Evidencia de contexto.} \id{EV-09} no origina requisitos de forma exclusiva. Se
        conserva en el catálogo por su valor de caracterización del proceso actual.
\end{enumerate}

% =====================================================================

\newpage
\section{Priorización y trazabilidad extendida}

\subsection{Clasificación de requisitos}

El documento especifica 75 elementos de requisito, distribuidos como sigue.

\begin{longtable}{|L{4.2cm}|C{2.0cm}|L{8.4cm}|}
\hline
\rowcolor{grisclaro}\textbf{Tipo} & \textbf{Cantidad} & \textbf{Identificadores} \\
\hline
\endhead
Requisito funcional (RF) & 39 & \id{RF-01} a \id{RF-39} \\ \hline
Requisito no funcional (RNF) & 18 & \id{RNF-01} a \id{RNF-18} \\ \hline
Restricción de diseño (RD) & 10 & \id{RD-01} a \id{RD-10} \\ \hline
Requisito legal (RL) & 8 & \id{RL-01} a \id{RL-08} \\ \hline
\rowcolor{grisclaro}\textbf{Total} & \textbf{75} & \\ \hline
\caption{Clasificación de los requisitos por tipo}
\end{longtable}

\noindent
La distribución MoSCoW de los requisitos funcionales es de 21 \emph{Must}, 16 \emph{Should} y 2
\emph{Could}. Los cuatro requisitos incorporados a partir de la hoja de liquidación semanal
(\id{RF-36} a \id{RF-39}) se clasifican como \emph{Must} los dos primeros ---el catálogo de tarifas
es prerrequisito del cálculo de remuneración y la liquidación es el cierre del ciclo semanal--- y
como \emph{Should} los dos restantes. No hay requisitos funcionales clasificados como \emph{Won't}: las exclusiones de
alcance se enuncian a nivel de producto en la §1.2.4 y se argumentan en la §5.3.4.

% ---------------------------------------------------------------------
\subsection{Historias de usuario de los requisitos \emph{Must}}

Conforme a la §3.3 de la guía de la Entrega 3, cada requisito funcional de prioridad \emph{Must}
cuenta con una historia de usuario en formato Connextra, con los criterios INVEST verificados y un
escenario de aceptación redactado en Gherkin. Se especifican 19 historias (\id{HU-01} a
\id{HU-19}), con sus criterios de aceptación (\id{CA-01} a \id{CA-19}).

\begin{longtable}{|C{1.3cm}|C{1.3cm}|L{12.4cm}|}
\hline
\rowcolor{grisclaro}\textbf{HU} & \textbf{RF} & \textbf{Historia de usuario (Connextra) y escenario de aceptación (Gherkin)} \\
\hline
\endhead

\id{HU-01} & \id{RF-01} & \textbf{Como} persona trabajadora del cultivo, \textbf{quiero} acceder al sistema con mis credenciales, \textbf{para} usar únicamente las funciones que corresponden a mi labor. \newline \emph{\id{CA-01}:} \textbf{Dado} que existe una cuenta con rol ``trabajador agrícola'', \textbf{cuando} inicia sesión, \textbf{entonces} el menú no presenta ninguna opción de administración \textbf{y} un intento de acceso directo a esa ruta es rechazado. \\ \hline

\id{HU-02} & \id{RF-02} & \textbf{Como} administrador general, \textbf{quiero} registrar los lotes con su variedad y superficie, \textbf{para} poder asociarles labores, diagnósticos y cosechas. \newline \emph{\id{CA-02}:} \textbf{Dado} un lote recién registrado con variedad asignada, \textbf{cuando} se abre el formulario de registro de labor, \textbf{entonces} el lote aparece como opción seleccionable. \\ \hline

\id{HU-03} & \id{RF-03} & \textbf{Como} administrador general, \textbf{quiero} organizar al personal por equipos funcionales, \textbf{para} asignar labores y evaluar el desempeño por cuadrilla. \newline \emph{\id{CA-03}:} \textbf{Dado} un trabajador asignado al equipo de polinización, \textbf{cuando} se consulta el equipo, \textbf{entonces} la persona figura en él \textbf{y} es asignable a labores de ese tipo. \\ \hline

\id{HU-04} & \id{RF-04} & \textbf{Como} trabajador agrícola, \textbf{quiero} registrar la labor que ejecuté en el lote, \textbf{para} que quede constancia sin depender de una libreta de papel. \newline \emph{\id{CA-04}:} \textbf{Dado} que ejecuté chapia en el lote 3, \textbf{cuando} registro la labor con su cantidad, \textbf{entonces} queda vinculada al lote, la fecha y mi identificación \textbf{y} aparece en el historial del lote. \\ \hline

\id{HU-05} & \id{RF-05} & \textbf{Como} técnico fitosanitario, \textbf{quiero} registrar el monitoreo sanitario del lote, \textbf{para} dejar trazabilidad de las revisiones y de las incidencias detectadas. \newline \emph{\id{CA-05}:} \textbf{Dado} un monitoreo con incidencia marcada, \textbf{cuando} se filtran los lotes con incidencia abierta, \textbf{entonces} ese lote aparece en el resultado. \\ \hline

\id{HU-06} & \id{RF-07} & \textbf{Como} asistente de administración, \textbf{quiero} fotografiar una palma afectada y obtener el diagnóstico, \textbf{para} identificar la plaga sin depender de conocimiento especializado. \newline \emph{\id{CA-06}:} \textbf{Dado} que capturo la imagen de una palma con afectación conocida, \textbf{cuando} solicito el análisis, \textbf{entonces} el sistema devuelve la clase, el nivel de confianza y una explicación de no más de 60 palabras \textbf{y} registra el diagnóstico con lote y fecha. \\ \hline

\id{HU-07} & \id{RF-08} & \textbf{Como} asesor técnico, \textbf{quiero} fotografiar una hoja y conocer la deficiencia nutricional, \textbf{para} decidir qué aplicar sin esperar el análisis de laboratorio. \newline \emph{\id{CA-07}:} \textbf{Dado} el lote con variedad asignada, \textbf{cuando} envío la imagen de una hoja con deficiencia conocida, \textbf{entonces} el sistema indica el elemento deficiente, propone una corrección \textbf{y} explica qué signo foliar la sustenta. \\ \hline

\id{HU-08} & \id{RF-10} & \textbf{Como} asesor técnico, \textbf{quiero} que el sistema conozca la variedad de cada lote, \textbf{para} que los umbrales de diagnóstico correspondan al material sembrado. \newline \emph{\id{CA-08}:} \textbf{Dado} un lote de variedad \emph{guineensis}, \textbf{cuando} se clasifica un racimo, \textbf{entonces} el umbral aplicado es de 3 frutos desprendidos; \textbf{y dado} un lote híbrido, \textbf{entonces} el umbral aplicado es de 5. \\ \hline

\id{HU-09} & \id{RF-12} & \textbf{Como} administrador general, \textbf{quiero} recibir alertas ante una afectación detectada, \textbf{para} intervenir antes de que se propague. \newline \emph{\id{CA-09}:} \textbf{Dado} un diagnóstico positivo registrado, \textbf{cuando} se completa el análisis, \textbf{entonces} se crea una alerta con lote, tipo, criticidad y estado abierto \textbf{y} resulta visible para el rol responsable. \\ \hline

\id{HU-10} & \id{RF-13} & \textbf{Como} jefe de polinización, \textbf{quiero} registrar el estado de la antesis y la verificación visual, \textbf{para} controlar la calidad de la labor que determina la producción. \newline \emph{\id{CA-10}:} \textbf{Dado} un lote en etapa de polinización, \textbf{cuando} se registra la jornada con estado de antesis, \textbf{entonces} el registro queda asociado al lote, la fecha y la persona \textbf{y} es consultable por fecha. \\ \hline

\id{HU-11} & \id{RF-14} & \textbf{Como} polinizador, \textbf{quiero} que el conteo de flores se calcule solo, \textbf{para} no perder rendimiento anotando planta por planta. \newline \emph{\id{CA-11}:} \textbf{Dado} que trabajé la jornada con rastreo activo, \textbf{cuando} cierro la jornada, \textbf{entonces} el sistema calcula el total de flores a partir de las marcaciones \textbf{y} no me solicita ningún registro manual por planta. \\ \hline

\id{HU-12} & \id{RF-18} & \textbf{Como} administrador general, \textbf{quiero} estimar la producción del período, \textbf{para} planificar la cosecha y la mano de obra necesaria. \newline \emph{\id{CA-12}:} \textbf{Dado} un censo de flores registrado y la variedad del lote, \textbf{cuando} solicito la estimación, \textbf{entonces} el sistema devuelve el tonelaje estimado del período \textbf{y} el supuesto de rendimiento empleado. \\ \hline

\id{HU-13} & \id{RF-19} & \textbf{Como} administrador general, \textbf{quiero} generar el consolidado semanal, \textbf{para} tomar decisiones sin recopilar reportes dispersos de mensajería. \newline \emph{\id{CA-13}:} \textbf{Dado} que existen registros en la semana, \textbf{cuando} genero el reporte, \textbf{entonces} obtengo labores, avance, producción y alertas del período \textbf{y} puedo exportarlo a PDF o a hoja de cálculo. \\ \hline

\id{HU-14} & \id{RF-21} & \textbf{Como} cosechador, \textbf{quiero} saber si un racimo está en punto de corte, \textbf{para} no cortar fruta verde que será penalizada. \newline \emph{\id{CA-14}:} \textbf{Dado} un lote híbrido, \textbf{cuando} fotografío la parte basal de un racimo con 4 frutos desprendidos, \textbf{entonces} el sistema lo clasifica como verde \textbf{y} explica que el umbral de la variedad es de 5. \\ \hline

\id{HU-15} & \id{RF-22} & \textbf{Como} administrador general, \textbf{quiero} conocer el porcentaje estimado de fruta verde antes de despachar, \textbf{para} evitar la penalización del lote en la planta extractora. \newline \emph{\id{CA-15}:} \textbf{Dado} un umbral de penalización del 3\%, \textbf{cuando} el porcentaje estimado del lote alcanza ese valor, \textbf{entonces} el sistema emite alerta preventiva \textbf{y} desaconseja el despacho. \\ \hline

\id{HU-16} & \id{RF-26} & \textbf{Como} administrador general, \textbf{quiero} registrar el presupuesto semanal de labores, \textbf{para} contrastar lo planificado con lo efectivamente ejecutado. \newline \emph{\id{CA-16}:} \textbf{Dado} un plan con meta de 1.123 coronas en el lote 3, \textbf{cuando} cierro la semana con 800 cumplidas, \textbf{entonces} el sistema reporta un cumplimiento del 71,2\% \textbf{y} un remanente de 323 coronas. \\ \hline

\id{HU-17} & \id{RF-28} & \textbf{Como} trabajador agrícola, \textbf{quiero} registrar mi avance en la unidad propia de cada labor, \textbf{para} que mi trabajo se contabilice como corresponde. \newline \emph{\id{CA-17}:} \textbf{Dado} que ejecuté 120 coronas, \textbf{cuando} registro el avance, \textbf{entonces} la cantidad se acumula a mi total semanal en la unidad ``coronas'' \textbf{y} el sistema rechaza una unidad incompatible con el tipo de labor. \\ \hline

\id{HU-18} & \id{RF-30} & \textbf{Como} asistente de administración, \textbf{quiero} reportar una afectación con foto y ubicación, \textbf{para} poder volver al punto exacto en mi siguiente visita. \newline \emph{\id{CA-18}:} \textbf{Dado} que reporto una incidencia con GPS activo, \textbf{cuando} la consulto días después, \textbf{entonces} recupero su fotografía y sus coordenadas \textbf{y} la incidencia figura en mi lista de pendientes de verificación. \\ \hline

\id{HU-19} & \id{RF-35} & \textbf{Como} jefe de campo, \textbf{quiero} registrar el avance de quien no tiene teléfono, \textbf{para} que ninguna persona quede sin registrar su trabajo. \newline \emph{\id{CA-19}:} \textbf{Dado} un trabajador sin dispositivo propio, \textbf{cuando} registro su avance de forma delegada, \textbf{entonces} el avance se atribuye a esa persona en su consolidado \textbf{y} la bitácora conserva mi identificación como registrador. \\ \hline
\caption{Historias de usuario y criterios de aceptación de los requisitos \emph{Must}}
\end{longtable}

\subsubsection{Verificación de los criterios INVEST}

Las diecinueve historias fueron evaluadas contra los seis criterios INVEST. La tabla siguiente
reporta el resultado y documenta las tres historias que requirieron ajuste.

\begin{longtable}{|L{2.4cm}|C{1.6cm}|L{11.0cm}|}
\hline
\rowcolor{grisclaro}\textbf{Criterio} & \textbf{Cumple} & \textbf{Observación} \\
\hline
\endhead
Independiente & 16 / 19 & \id{HU-11} depende de \id{HU-10}, \id{HU-15} de \id{HU-14} y \id{HU-17} de \id{HU-04}. Se documenta la dependencia en la matriz en lugar de forzar la independencia, dado que responde a la secuencia real del proceso. \\ \hline
Negociable & 19 / 19 & Ninguna historia prescribe solución técnica; los umbrales son parametrizables (\id{RD-08}). \\ \hline
Valiosa & 19 / 19 & Cada historia declara el beneficio esperado y se traza a evidencia de campo. \\ \hline
Estimable & 19 / 19 & Todas admiten estimación en la escala Fibonacci empleada en el cálculo WSJF de la §5.3.3. \\ \hline
Pequeña & 17 / 19 & \id{HU-06} y \id{HU-07} superan el tamaño recomendado por incluir el componente de IA; se sugiere descomponerlas en captura, inferencia y explicación durante la planificación de la implementación. \\ \hline
Verificable & 19 / 19 & Todas cuentan con escenario Gherkin con criterio observable. \\ \hline
\caption{Verificación de los criterios INVEST sobre las historias de usuario}
\end{longtable}

% ---------------------------------------------------------------------
\subsection{Priorización combinada}

La priorización emplea tres instrumentos complementarios: MoSCoW para el compromiso de alcance, el
modelo de Kano para la relación con la satisfacción del cliente, y WSJF para el orden de
implementación.

\subsubsection{Justificación de la clasificación MoSCoW}

\textbf{\emph{Must} (19 RF).} Corresponden a tres condiciones. Primero, la \textbf{base operativa
indispensable}: sin autenticación por rol (\id{RF-01}), estructura de lotes (\id{RF-02}), personal
(\id{RF-03}) y registro de labor y avance (\id{RF-04}, \id{RF-28}, \id{RF-35}) el sistema no
sustituye al proceso manual que pretende reemplazar. Segundo, el \textbf{diferencial de valor}: el
análisis por imagen (\id{RF-07}, \id{RF-08}, \id{RF-21}) y las alertas (\id{RF-12}) son las
funciones mejor valoradas por el personal en el cuestionario ($4{,}31$ y $4{,}27$ sobre 5,
\id{EV-12}) y las señaladas de forma independiente por la administración (\id{EV-01}) y la asesoría
técnica (\id{EV-02}). Tercero, la \textbf{consecuencia económica directa}: la polinización determina
la producción de la finca (\id{EV-03}), su conteo condiciona la evaluación del trabajo
(\id{RF-13}, \id{RF-14}), y el porcentaje de fruta verde determina la penalización del lote en la
planta extractora (\id{RF-22}, \id{EV-04}).

\textbf{\emph{Should} (14 RF).} Aportan valor sustantivo pero el sistema entrega su propuesta de
valor sin ellos en la primera versión: seguimiento de etapas, ficha de plaga, análisis de
suelo, recorridos, evaluación de desempeño, clima, historial, registro del ticket de la extractora,
trazabilidad de la penalización, ventana corte-entrega, arrastre semanal, consulta de remuneración,
lista de pendientes y exportación hacia la extractora.

\textbf{\emph{Could} (2 RF).} \id{RF-33} (guía de remisión) y \id{RF-34} (solicitud de insumos) son
mejoras deseables de bajo acoplamiento con el núcleo.

\textbf{\emph{Won't} en esta versión.} Se excluyen explícitamente y con argumento: la
\textbf{gestión de nómina y pago}, porque el sistema calcula el avance pero la liquidación
involucra información laboral y tributaria cuyo tratamiento excede el consentimiento obtenido
(\id{RD-04}, \id{RL-02}); la \textbf{facturación y contabilidad}, por no corresponder al manejo del
cultivo; la \textbf{comercialización de la producción}, que se negocia fuera del sistema; la
\textbf{integración directa con el sistema de la planta extractora}, porque pertenece a una
organización externa con sistema propio y su interfaz no está disponible (\id{EV-04}); y la
\textbf{operación multi-finca}, que el cliente administra pero que no forma parte del alcance
acordado para la v1 (§1.2.2).

\subsubsection{Modelo de Kano}

La clasificación se sustenta en la distancia entre la valoración declarada en el cuestionario y la
expectativa implícita del proceso actual.

\begin{longtable}{|L{2.8cm}|L{5.4cm}|L{7.0cm}|}
\hline
\rowcolor{grisclaro}\textbf{Categoría} & \textbf{Requisitos} & \textbf{Fundamento} \\
\hline
\endhead
Básica & \id{RF-01}, \id{RF-02}, \id{RF-03}, \id{RF-04} & Su ausencia genera rechazo, su presencia no genera entusiasmo. El registro de labor ya existe en papel: el sistema debe igualarlo, no impresiona por hacerlo. \\ \hline
De desempeño & \id{RF-19}, \id{RF-26}, \id{RF-28}, \id{RF-29}, \id{RF-13} & La satisfacción crece de forma proporcional a la calidad de ejecución. Valoración intermedia en el cuestionario ($3{,}95$ para el registro de labores). \\ \hline
\rowcolor{grisclaro}De deleite & \id{RF-07}, \id{RF-08}, \id{RF-12}, \id{RF-21} & Generan entusiasmo desproporcionado respecto de la expectativa. Son las funciones mejor valoradas ($4{,}27$ y $4{,}31$) y las únicas mencionadas de forma espontánea por los cinco perfiles entrevistados. \\ \hline
Indiferente & \id{RF-33}, \id{RF-34} & No inciden en la satisfacción percibida por la persona usuaria final. \\ \hline
De rechazo potencial & \id{RF-15} & El rastreo de recorridos puede generar insatisfacción: el 25{,}8\% del personal manifiesta reservas (\id{EV-12}). Se mitiga con consentimiento revocable (\id{RL-03}). \\ \hline
\caption{Clasificación de los requisitos según el modelo de Kano}
\end{longtable}

\noindent
La clasificación de \id{RF-15} como \emph{requisito de rechazo potencial} es deliberada y
constituye un hallazgo. La funcionalidad fue solicitada por la administración (\id{EV-01},
\id{EV-02}) y por la jefatura de polinización (\id{EV-03}), pero una de cada cuatro personas
destinatarias expresa reservas sobre ella. Kano permite representar esta tensión: un requisito puede
ser valioso para quien lo solicita y adverso para quien lo padece. La mitigación no es técnica sino
de diseño de la relación: consentimiento específico, revocable y sin consecuencia laboral.

\subsubsection{Valor de negocio mediante WSJF}

Se aplica el método del Trabajo Ponderado Más Corto Primero:
\[
\mathrm{WSJF} = \frac{\text{Valor de negocio} + \text{Criticidad temporal} + \text{Reducción de riesgo}}{\text{Tamaño del trabajo}}
\]
Los cuatro componentes se puntúan en escala Fibonacci (1, 2, 3, 5, 8, 13). La tabla presenta los
diez requisitos evaluados de mayor valor de negocio, ordenados de forma descendente por WSJF.

\begin{longtable}{|C{1.4cm}|C{1.3cm}|C{1.3cm}|C{1.3cm}|C{1.3cm}|C{1.4cm}|L{5.4cm}|}
\hline
\rowcolor{grisclaro}\textbf{RF} & \textbf{Valor} & \textbf{Crit. temp.} & \textbf{Red. riesgo} & \textbf{Tamaño} & \textbf{WSJF} & \textbf{Justificación} \\
\hline
\endhead
\id{RF-35} & 5 & 8 & 8 & 3 & \textbf{7,00} & Bajo costo y alto impacto: sin él, el 11,3\% del personal queda excluido del sistema (\id{EV-12}). \\ \hline
\id{RF-21} & 13 & 13 & 8 & 8 & \textbf{4,25} & Evita la penalización económica del lote; la temporada de baja producción (jul.--oct.) es cuando más fruta verde se corta (\id{EV-04}). \\ \hline
\id{RF-22} & 13 & 13 & 8 & 8 & \textbf{4,25} & Actúa antes del despacho, momento en que la pérdida aún es evitable. \\ \hline
\id{RF-04} & 8 & 8 & 5 & 5 & \textbf{4,20} & Base de todo el registro operativo; sin él no hay datos para ninguna otra función. \\ \hline
\id{RF-12} & 8 & 8 & 5 & 5 & \textbf{4,20} & Función mejor valorada del cuestionario (4,31/5). \\ \hline
\id{RF-01} & 5 & 8 & 8 & 5 & \textbf{4,20} & Precondición técnica y legal de todo lo demás (\id{RL-01}). \\ \hline
\id{RF-07} & 13 & 8 & 8 & 8 & \textbf{3,62} & Mayor dificultad declarada por el personal (40,3\%, \id{EV-12}). \\ \hline
\id{RF-28} & 8 & 5 & 5 & 5 & \textbf{3,60} & Habilita el cálculo de cumplimiento y de remuneración por avance. \\ \hline
\id{RF-14} & 13 & 8 & 5 & 8 & \textbf{3,25} & La polinización determina la producción; el conteo manual afecta el rendimiento (\id{EV-03}). \\ \hline
\id{RF-08} & 8 & 5 & 5 & 8 & \textbf{2,25} & Alto valor técnico, pero el diagnóstico nutricional admite mayor demora que el sanitario. \\ \hline

\caption{Valor de negocio mediante WSJF (diez requisitos de mayor prioridad)}
\end{longtable}

\noindent
\textbf{Lectura del resultado.} \id{RF-35} y \id{RF-10} encabezan el orden de implementación con un
WSJF de 7,00 cada uno,
pese a no ser el requisito de mayor valor de negocio. La razón es su tamaño reducido frente a un
impacto alto: se trata de una variante del formulario de registro que evita excluir a una de cada
nueve personas destinatarias. \id{RF-10} alcanza la misma posición por una razón distinta: es
prerrequisito de \id{RF-21}, dado que el umbral de desprendimiento se obtiene de la variedad del
lote. El método revela así una dependencia técnica que ni MoSCoW ni Kano hacen visible.

Ambos casos ilustran la utilidad del método frente a una priorización basada únicamente en valor
percibido, que los habría situado muy por debajo. La lectura inversa también es informativa:
\id{RF-07} y \id{RF-08}, las dos funciones mejor valoradas por las personas usuarias, obtienen WSJF
de 3,62 y 2,25 por su elevado tamaño de trabajo. Esto no altera su condición de \emph{Must}, pero
indica que el desarrollo debe iniciarse por la base operativa mientras los componentes de
inteligencia artificial se construyen en paralelo. El archivo
\id{04\_Trazabilidad/}\allowbreak\id{priorizacion\_moscow\_kano.csv} contiene la tabla completa para los 35
requisitos funcionales.

% ---------------------------------------------------------------------
\subsection{Matriz de trazabilidad extendida}

La matriz enlaza, para cada elemento, la cadena completa
Ley $\rightarrow$ Objetivo $\rightarrow$ Interesado $\rightarrow$ Evidencia $\rightarrow$
Requisito $\rightarrow$ Caso de uso $\rightarrow$ Historia de usuario $\rightarrow$
Criterio de aceptación $\rightarrow$ Componente $\rightarrow$ Mockup.

La versión completa, con las trece columnas exigidas por la §8.8 de la guía, reside en
\id{04\_Trazabilidad/}\allowbreak\id{matriz\_trazabilidad.csv} y contiene \textbf{52 filas}. La tabla siguiente
reproduce una vista reducida con las columnas de mayor densidad informativa; se conservan las 48
filas para permitir la verificación de cobertura. Las cuatro últimas corresponden a los requisitos derivados de la hoja de liquidación semanal (\id{EV-13}).

\begin{longtable}{|C{0.9cm}|C{1.5cm}|C{1.6cm}|C{1.5cm}|C{1.3cm}|C{1.3cm}|C{1.3cm}|C{1.5cm}|}
\hline
\rowcolor{grisclaro}\textbf{\#} & \textbf{Interesado} & \textbf{ID-EV} & \textbf{ID-RF} & \textbf{Tipo} & \textbf{ID-CU} & \textbf{ID-HU} & \textbf{ID-Mockup} \\
\hline
\endhead
1 & Administrador & EV-01, EV-09 & RF-01 & RF & CU-11 & HU-01 & MU-01 \\ \hline
2 & Administrador & EV-01, EV-02 & RF-02 & RF & CU-12 & HU-02 & MU-05 \\ \hline
3 & Administrador & EV-01, EV-11 & RF-03 & RF & CU-13 & HU-03 & --- \\ \hline
4 & Trabajador & EV-05, EV-06 & RF-04 & RF & CU-04 & HU-04 & MU-04 \\ \hline
5 & Trabajador & EV-11 & RF-04 & RF & CU-04 & HU-04 & MU-04 \\ \hline
6 & Técnico fitosan. & EV-01, EV-02 & RF-05 & RF & CU-15 & HU-05 & MU-05 \\ \hline
7 & Administrador & EV-01 & RF-06 & RF & CU-12 & --- & MU-05 \\ \hline
8 & Asistente admin. & EV-01, EV-02 & RF-07 & RF & CU-01 & HU-06 & MU-03 \\ \hline
9 & Asistente admin. & EV-07 & RF-07 & RF & CU-01 & HU-06 & MU-03 \\ \hline
10 & Asesor técnico & EV-02 & RF-08 & RF & CU-02 & HU-07 & MU-03 \\ \hline
11 & Asesor técnico & EV-04 & RF-08 & RF & CU-02 & HU-07 & MU-03 \\ \hline
12 & Técnico fitosan. & EV-01, EV-02 & RF-09 & RF & CU-01 & --- & MU-03 \\ \hline
13 & Asesor técnico & EV-04 & RF-10 & RF & CU-12 & HU-08 & MU-05 \\ \hline
14 & Asesor técnico & EV-02 & RF-11 & RF & CU-02 & --- & MU-05 \\ \hline
15 & Administrador & EV-01, EV-02 & RF-12 & RF & CU-16 & HU-09 & MU-08 \\ \hline
16 & Trabajador & EV-12 & RF-12 & RF & CU-16 & HU-09 & MU-08 \\ \hline
17 & Jefe polinización & EV-03 & RF-13 & RF & CU-03 & HU-10 & MU-04 \\ \hline
18 & Jefe polinización & EV-03 & RF-14 & RF & CU-03 & HU-11 & MU-06 \\ \hline
19 & Administrador & EV-01, EV-03 & RF-15 & RF & CU-03 & --- & MU-06 \\ \hline
20 & Administrador & EV-01, EV-02 & RF-16 & RF & CU-05 & --- & MU-07 \\ \hline
21 & Administrador & EV-01, EV-02 & RF-17 & RF & CU-05 & --- & MU-05 \\ \hline
22 & Administrador & EV-01, EV-02 & RF-18 & RF & CU-17 & HU-12 & MU-07 \\ \hline
23 & Administrador & EV-01, EV-11 & RF-19 & RF & CU-05 & HU-13 & MU-07 \\ \hline
24 & Administrador & EV-01 & RF-20 & RF & CU-05 & --- & MU-05 \\ \hline
25 & Téc. extractora & EV-04 & RF-21 & RF & CU-06 & HU-14 & MU-03 \\ \hline
26 & Téc. extractora & EV-08 & RF-21 & RF & CU-06 & HU-14 & MU-03 \\ \hline
27 & Administrador & EV-04, EV-08 & RF-22 & RF & CU-06 & HU-15 & MU-08 \\ \hline
28 & Administrador & EV-04 & RF-23 & RF & CU-10 & --- & MU-07 \\ \hline
29 & Téc. extractora & EV-04, EV-08 & RF-24 & RF & CU-10 & --- & MU-07 \\ \hline
30 & Téc. extractora & EV-04 & RF-25 & RF & CU-10 & --- & MU-07 \\ \hline
31 & Administrador & EV-11 & RF-26 & RF & CU-07 & HU-16 & MU-02 \\ \hline
32 & Administrador & EV-11 & RF-27 & RF & CU-07 & --- & MU-02 \\ \hline
33 & Trabajador & EV-05, EV-11 & RF-28 & RF & CU-04 & HU-17 & MU-04 \\ \hline
34 & Trabajador & EV-06, EV-07 & RF-29 & RF & CU-08 & --- & MU-02 \\ \hline
35 & Asistente admin. & EV-07 & RF-30 & RF & CU-09 & HU-18 & MU-06 \\ \hline
36 & Asistente admin. & EV-07 & RF-31 & RF & CU-09 & --- & MU-08 \\ \hline
37 & Téc. extractora & EV-04 & RF-32 & RF & CU-17 & --- & MU-07 \\ \hline
38 & Téc. extractora & EV-04 & RF-33 & RF & CU-10 & --- & MU-07 \\ \hline
39 & Administrador & EV-11 & RF-34 & RF & CU-07 & --- & MU-02 \\ \hline
40 & Jefe de campo & EV-12, EV-05 & RF-35 & RF & CU-04 & HU-19 & MU-04 \\ \hline
41 & Téc. extractora & EV-04 & RNF-01 & RNF & CU-06 & HU-14 & MU-03 \\ \hline
42 & Trabajador & EV-07, EV-12 & RNF-11 & RNF & CU-04 & HU-04 & MU-04 \\ \hline
43 & Jefe polinización & EV-03 & RNF-05 & RNF & CU-03 & HU-11 & MU-06 \\ \hline
44 & Trabajador & EV-12 & RNF-08 & RNF & CU-04 & HU-04 & MU-04 \\ \hline
45 & Trabajador & EV-12 & RNF-14 & RNF & CU-03 & HU-11 & MU-06 \\ \hline
46 & Asistente admin. & EV-01, EV-07 & RNF-16 & RNF & CU-18 & HU-06 & MU-03 \\ \hline
47 & Trabajador & EV-12 & RD-10 & RD & CU-04 & HU-19 & MU-04 \\ \hline
48 & Téc. extractora & EV-08 & RD-07 & RD & CU-06 & HU-14 & MU-03 \\ \hline
49 & Administrador & EV-13 & RF-36 & RF & CU-08 & --- & MU-02 \\ \hline
50 & Administrador & EV-13, EV-11 & RF-37 & RF & CU-07 & --- & MU-02 \\ \hline
51 & Administrador & EV-13 & RF-38 & RF & CU-07 & --- & MU-02 \\ \hline
52 & Administrador & EV-13 & RF-39 & RF & CU-07 & --- & MU-02 \\ \hline
\caption{Matriz de trazabilidad extendida (vista reducida, 52 filas)}
\end{longtable}

\subsubsection{Diagnóstico de cobertura}

\begin{longtable}{|L{5.6cm}|C{2.4cm}|L{6.8cm}|}
\hline
\rowcolor{grisclaro}\textbf{Indicador} & \textbf{Resultado} & \textbf{Interpretación} \\
\hline
\endhead
Requisitos funcionales trazados a evidencia & 39 / 39 (100\%) & Ningún requisito carece de fuente verificable. \\ \hline
Requisitos \emph{Must} con caso de uso & 19 / 19 (100\%) & Se corrige la deficiencia de la Entrega 2, donde \id{RF-17} y \id{RF-20} carecían de caso de uso. \\ \hline
Requisitos \emph{Must} con historia y criterio & 19 / 19 (100\%) & \\ \hline
Evidencias que originan al menos un requisito & 12 / 13 (92,3\%) & \id{EV-09} (observación de campo) aporta contexto y respalda requisitos junto con otras fuentes, sin originar ninguno de forma exclusiva. \\ \hline
Requisitos sustentados en una sola evidencia & 9 / 39 (23,1\%) & Véase el análisis de elementos huérfanos. \\ \hline
Requisitos con mockup asociado & 36 / 39 (92,3\%) & \id{RF-03}, \id{RF-27} y \id{RF-34} operan sobre pantallas no prototipadas en esta entrega. \\ \hline
\caption{Diagnóstico de cobertura de la matriz de trazabilidad}
\end{longtable}

\subsubsection{Elementos huérfanos y vacíos de cobertura}

Se declaran de forma explícita las debilidades detectadas, conforme al criterio de completitud de la
guía.

\begin{enumerate}[nosep]
  \item \textbf{Requisitos con fuente única.} \id{RF-11}, \id{RF-23}, \id{RF-25}, \id{RF-32},
        \id{RF-33} y \id{RF-34} se sustentan en una sola evidencia. Los cuatro derivados de
        \id{EV-04} conciernen a la relación con la planta extractora y no pudieron triangularse con
        una segunda fuente independiente de esa organización, dado que \id{EV-08} corresponde a la
        misma sesión de campo. Se declara como limitación.
  \item \textbf{Requisitos sin prototipo de interfaz.} \id{RF-03}, \id{RF-27} y \id{RF-34}. Los tres
        son de prioridad no \emph{Must}, por lo que la ausencia no incumple la exigencia de
        prototipar las pantallas \emph{Must}.
  \item \textbf{Ausencia de \emph{Won't} funcionales.} Las exclusiones se enuncian a nivel de
        producto y no como requisitos individuales descartados. Se argumentan en la §5.3.1.
  \item \textbf{Evidencia de contexto.} \id{EV-09} no origina requisitos de forma exclusiva. Se
        conserva en el catálogo por su valor de caracterización del proceso actual.
\end{enumerate}

% =====================================================================

\newpage
\section{Priorización y trazabilidad extendida}

\subsection{Clasificación de requisitos}

El documento especifica 75 elementos de requisito, distribuidos como sigue.

\begin{longtable}{|L{4.2cm}|C{2.0cm}|L{8.4cm}|}
\hline
\rowcolor{grisclaro}\textbf{Tipo} & \textbf{Cantidad} & \textbf{Identificadores} \\
\hline
\endhead
Requisito funcional (RF) & 39 & \id{RF-01} a \id{RF-39} \\ \hline
Requisito no funcional (RNF) & 18 & \id{RNF-01} a \id{RNF-18} \\ \hline
Restricción de diseño (RD) & 10 & \id{RD-01} a \id{RD-10} \\ \hline
Requisito legal (RL) & 8 & \id{RL-01} a \id{RL-08} \\ \hline
\rowcolor{grisclaro}\textbf{Total} & \textbf{75} & \\ \hline
\caption{Clasificación de los requisitos por tipo}
\end{longtable}

\noindent
La distribución MoSCoW de los requisitos funcionales es de 21 \emph{Must}, 16 \emph{Should} y 2
\emph{Could}. Los cuatro requisitos incorporados a partir de la hoja de liquidación semanal
(\id{RF-36} a \id{RF-39}) se clasifican como \emph{Must} los dos primeros ---el catálogo de tarifas
es prerrequisito del cálculo de remuneración y la liquidación es el cierre del ciclo semanal--- y
como \emph{Should} los dos restantes. No hay requisitos funcionales clasificados como \emph{Won't}: las exclusiones de
alcance se enuncian a nivel de producto en la §1.2.4 y se argumentan en la §5.3.4.

% ---------------------------------------------------------------------
\subsection{Historias de usuario de los requisitos \emph{Must}}

Conforme a la §3.3 de la guía de la Entrega 3, cada requisito funcional de prioridad \emph{Must}
cuenta con una historia de usuario en formato Connextra, con los criterios INVEST verificados y un
escenario de aceptación redactado en Gherkin. Se especifican 19 historias (\id{HU-01} a
\id{HU-19}), con sus criterios de aceptación (\id{CA-01} a \id{CA-19}).

\begin{longtable}{|C{1.3cm}|C{1.3cm}|L{12.4cm}|}
\hline
\rowcolor{grisclaro}\textbf{HU} & \textbf{RF} & \textbf{Historia de usuario (Connextra) y escenario de aceptación (Gherkin)} \\
\hline
\endhead

\id{HU-01} & \id{RF-01} & \textbf{Como} persona trabajadora del cultivo, \textbf{quiero} acceder al sistema con mis credenciales, \textbf{para} usar únicamente las funciones que corresponden a mi labor. \newline \emph{\id{CA-01}:} \textbf{Dado} que existe una cuenta con rol ``trabajador agrícola'', \textbf{cuando} inicia sesión, \textbf{entonces} el menú no presenta ninguna opción de administración \textbf{y} un intento de acceso directo a esa ruta es rechazado. \\ \hline

\id{HU-02} & \id{RF-02} & \textbf{Como} administrador general, \textbf{quiero} registrar los lotes con su variedad y superficie, \textbf{para} poder asociarles labores, diagnósticos y cosechas. \newline \emph{\id{CA-02}:} \textbf{Dado} un lote recién registrado con variedad asignada, \textbf{cuando} se abre el formulario de registro de labor, \textbf{entonces} el lote aparece como opción seleccionable. \\ \hline

\id{HU-03} & \id{RF-03} & \textbf{Como} administrador general, \textbf{quiero} organizar al personal por equipos funcionales, \textbf{para} asignar labores y evaluar el desempeño por cuadrilla. \newline \emph{\id{CA-03}:} \textbf{Dado} un trabajador asignado al equipo de polinización, \textbf{cuando} se consulta el equipo, \textbf{entonces} la persona figura en él \textbf{y} es asignable a labores de ese tipo. \\ \hline

\id{HU-04} & \id{RF-04} & \textbf{Como} trabajador agrícola, \textbf{quiero} registrar la labor que ejecuté en el lote, \textbf{para} que quede constancia sin depender de una libreta de papel. \newline \emph{\id{CA-04}:} \textbf{Dado} que ejecuté chapia en el lote 3, \textbf{cuando} registro la labor con su cantidad, \textbf{entonces} queda vinculada al lote, la fecha y mi identificación \textbf{y} aparece en el historial del lote. \\ \hline

\id{HU-05} & \id{RF-05} & \textbf{Como} técnico fitosanitario, \textbf{quiero} registrar el monitoreo sanitario del lote, \textbf{para} dejar trazabilidad de las revisiones y de las incidencias detectadas. \newline \emph{\id{CA-05}:} \textbf{Dado} un monitoreo con incidencia marcada, \textbf{cuando} se filtran los lotes con incidencia abierta, \textbf{entonces} ese lote aparece en el resultado. \\ \hline

\id{HU-06} & \id{RF-07} & \textbf{Como} asistente de administración, \textbf{quiero} fotografiar una palma afectada y obtener el diagnóstico, \textbf{para} identificar la plaga sin depender de conocimiento especializado. \newline \emph{\id{CA-06}:} \textbf{Dado} que capturo la imagen de una palma con afectación conocida, \textbf{cuando} solicito el análisis, \textbf{entonces} el sistema devuelve la clase, el nivel de confianza y una explicación de no más de 60 palabras \textbf{y} registra el diagnóstico con lote y fecha. \\ \hline

\id{HU-07} & \id{RF-08} & \textbf{Como} asesor técnico, \textbf{quiero} fotografiar una hoja y conocer la deficiencia nutricional, \textbf{para} decidir qué aplicar sin esperar el análisis de laboratorio. \newline \emph{\id{CA-07}:} \textbf{Dado} el lote con variedad asignada, \textbf{cuando} envío la imagen de una hoja con deficiencia conocida, \textbf{entonces} el sistema indica el elemento deficiente, propone una corrección \textbf{y} explica qué signo foliar la sustenta. \\ \hline

\id{HU-08} & \id{RF-10} & \textbf{Como} asesor técnico, \textbf{quiero} que el sistema conozca la variedad de cada lote, \textbf{para} que los umbrales de diagnóstico correspondan al material sembrado. \newline \emph{\id{CA-08}:} \textbf{Dado} un lote de variedad \emph{guineensis}, \textbf{cuando} se clasifica un racimo, \textbf{entonces} el umbral aplicado es de 3 frutos desprendidos; \textbf{y dado} un lote híbrido, \textbf{entonces} el umbral aplicado es de 5. \\ \hline

\id{HU-09} & \id{RF-12} & \textbf{Como} administrador general, \textbf{quiero} recibir alertas ante una afectación detectada, \textbf{para} intervenir antes de que se propague. \newline \emph{\id{CA-09}:} \textbf{Dado} un diagnóstico positivo registrado, \textbf{cuando} se completa el análisis, \textbf{entonces} se crea una alerta con lote, tipo, criticidad y estado abierto \textbf{y} resulta visible para el rol responsable. \\ \hline

\id{HU-10} & \id{RF-13} & \textbf{Como} jefe de polinización, \textbf{quiero} registrar el estado de la antesis y la verificación visual, \textbf{para} controlar la calidad de la labor que determina la producción. \newline \emph{\id{CA-10}:} \textbf{Dado} un lote en etapa de polinización, \textbf{cuando} se registra la jornada con estado de antesis, \textbf{entonces} el registro queda asociado al lote, la fecha y la persona \textbf{y} es consultable por fecha. \\ \hline

\id{HU-11} & \id{RF-14} & \textbf{Como} polinizador, \textbf{quiero} que el conteo de flores se calcule solo, \textbf{para} no perder rendimiento anotando planta por planta. \newline \emph{\id{CA-11}:} \textbf{Dado} que trabajé la jornada con rastreo activo, \textbf{cuando} cierro la jornada, \textbf{entonces} el sistema calcula el total de flores a partir de las marcaciones \textbf{y} no me solicita ningún registro manual por planta. \\ \hline

\id{HU-12} & \id{RF-18} & \textbf{Como} administrador general, \textbf{quiero} estimar la producción del período, \textbf{para} planificar la cosecha y la mano de obra necesaria. \newline \emph{\id{CA-12}:} \textbf{Dado} un censo de flores registrado y la variedad del lote, \textbf{cuando} solicito la estimación, \textbf{entonces} el sistema devuelve el tonelaje estimado del período \textbf{y} el supuesto de rendimiento empleado. \\ \hline

\id{HU-13} & \id{RF-19} & \textbf{Como} administrador general, \textbf{quiero} generar el consolidado semanal, \textbf{para} tomar decisiones sin recopilar reportes dispersos de mensajería. \newline \emph{\id{CA-13}:} \textbf{Dado} que existen registros en la semana, \textbf{cuando} genero el reporte, \textbf{entonces} obtengo labores, avance, producción y alertas del período \textbf{y} puedo exportarlo a PDF o a hoja de cálculo. \\ \hline

\id{HU-14} & \id{RF-21} & \textbf{Como} cosechador, \textbf{quiero} saber si un racimo está en punto de corte, \textbf{para} no cortar fruta verde que será penalizada. \newline \emph{\id{CA-14}:} \textbf{Dado} un lote híbrido, \textbf{cuando} fotografío la parte basal de un racimo con 4 frutos desprendidos, \textbf{entonces} el sistema lo clasifica como verde \textbf{y} explica que el umbral de la variedad es de 5. \\ \hline

\id{HU-15} & \id{RF-22} & \textbf{Como} administrador general, \textbf{quiero} conocer el porcentaje estimado de fruta verde antes de despachar, \textbf{para} evitar la penalización del lote en la planta extractora. \newline \emph{\id{CA-15}:} \textbf{Dado} un umbral de penalización del 3\%, \textbf{cuando} el porcentaje estimado del lote alcanza ese valor, \textbf{entonces} el sistema emite alerta preventiva \textbf{y} desaconseja el despacho. \\ \hline

\id{HU-16} & \id{RF-26} & \textbf{Como} administrador general, \textbf{quiero} registrar el presupuesto semanal de labores, \textbf{para} contrastar lo planificado con lo efectivamente ejecutado. \newline \emph{\id{CA-16}:} \textbf{Dado} un plan con meta de 1.123 coronas en el lote 3, \textbf{cuando} cierro la semana con 800 cumplidas, \textbf{entonces} el sistema reporta un cumplimiento del 71,2\% \textbf{y} un remanente de 323 coronas. \\ \hline

\id{HU-17} & \id{RF-28} & \textbf{Como} trabajador agrícola, \textbf{quiero} registrar mi avance en la unidad propia de cada labor, \textbf{para} que mi trabajo se contabilice como corresponde. \newline \emph{\id{CA-17}:} \textbf{Dado} que ejecuté 120 coronas, \textbf{cuando} registro el avance, \textbf{entonces} la cantidad se acumula a mi total semanal en la unidad ``coronas'' \textbf{y} el sistema rechaza una unidad incompatible con el tipo de labor. \\ \hline

\id{HU-18} & \id{RF-30} & \textbf{Como} asistente de administración, \textbf{quiero} reportar una afectación con foto y ubicación, \textbf{para} poder volver al punto exacto en mi siguiente visita. \newline \emph{\id{CA-18}:} \textbf{Dado} que reporto una incidencia con GPS activo, \textbf{cuando} la consulto días después, \textbf{entonces} recupero su fotografía y sus coordenadas \textbf{y} la incidencia figura en mi lista de pendientes de verificación. \\ \hline

\id{HU-19} & \id{RF-35} & \textbf{Como} jefe de campo, \textbf{quiero} registrar el avance de quien no tiene teléfono, \textbf{para} que ninguna persona quede sin registrar su trabajo. \newline \emph{\id{CA-19}:} \textbf{Dado} un trabajador sin dispositivo propio, \textbf{cuando} registro su avance de forma delegada, \textbf{entonces} el avance se atribuye a esa persona en su consolidado \textbf{y} la bitácora conserva mi identificación como registrador. \\ \hline
\caption{Historias de usuario y criterios de aceptación de los requisitos \emph{Must}}
\end{longtable}

\subsubsection{Verificación de los criterios INVEST}

Las diecinueve historias fueron evaluadas contra los seis criterios INVEST. La tabla siguiente
reporta el resultado y documenta las tres historias que requirieron ajuste.

\begin{longtable}{|L{2.4cm}|C{1.6cm}|L{11.0cm}|}
\hline
\rowcolor{grisclaro}\textbf{Criterio} & \textbf{Cumple} & \textbf{Observación} \\
\hline
\endhead
Independiente & 16 / 19 & \id{HU-11} depende de \id{HU-10}, \id{HU-15} de \id{HU-14} y \id{HU-17} de \id{HU-04}. Se documenta la dependencia en la matriz en lugar de forzar la independencia, dado que responde a la secuencia real del proceso. \\ \hline
Negociable & 19 / 19 & Ninguna historia prescribe solución técnica; los umbrales son parametrizables (\id{RD-08}). \\ \hline
Valiosa & 19 / 19 & Cada historia declara el beneficio esperado y se traza a evidencia de campo. \\ \hline
Estimable & 19 / 19 & Todas admiten estimación en la escala Fibonacci empleada en el cálculo WSJF de la §5.3.3. \\ \hline
Pequeña & 17 / 19 & \id{HU-06} y \id{HU-07} superan el tamaño recomendado por incluir el componente de IA; se sugiere descomponerlas en captura, inferencia y explicación durante la planificación de la implementación. \\ \hline
Verificable & 19 / 19 & Todas cuentan con escenario Gherkin con criterio observable. \\ \hline
\caption{Verificación de los criterios INVEST sobre las historias de usuario}
\end{longtable}

% ---------------------------------------------------------------------
\subsection{Priorización combinada}

La priorización emplea tres instrumentos complementarios: MoSCoW para el compromiso de alcance, el
modelo de Kano para la relación con la satisfacción del cliente, y WSJF para el orden de
implementación.

\subsubsection{Justificación de la clasificación MoSCoW}

\textbf{\emph{Must} (19 RF).} Corresponden a tres condiciones. Primero, la \textbf{base operativa
indispensable}: sin autenticación por rol (\id{RF-01}), estructura de lotes (\id{RF-02}), personal
(\id{RF-03}) y registro de labor y avance (\id{RF-04}, \id{RF-28}, \id{RF-35}) el sistema no
sustituye al proceso manual que pretende reemplazar. Segundo, el \textbf{diferencial de valor}: el
análisis por imagen (\id{RF-07}, \id{RF-08}, \id{RF-21}) y las alertas (\id{RF-12}) son las
funciones mejor valoradas por el personal en el cuestionario ($4{,}31$ y $4{,}27$ sobre 5,
\id{EV-12}) y las señaladas de forma independiente por la administración (\id{EV-01}) y la asesoría
técnica (\id{EV-02}). Tercero, la \textbf{consecuencia económica directa}: la polinización determina
la producción de la finca (\id{EV-03}), su conteo condiciona la evaluación del trabajo
(\id{RF-13}, \id{RF-14}), y el porcentaje de fruta verde determina la penalización del lote en la
planta extractora (\id{RF-22}, \id{EV-04}).

\textbf{\emph{Should} (14 RF).} Aportan valor sustantivo pero el sistema entrega su propuesta de
valor sin ellos en la primera versión: seguimiento de etapas, ficha de plaga, análisis de
suelo, recorridos, evaluación de desempeño, clima, historial, registro del ticket de la extractora,
trazabilidad de la penalización, ventana corte-entrega, arrastre semanal, consulta de remuneración,
lista de pendientes y exportación hacia la extractora.

\textbf{\emph{Could} (2 RF).} \id{RF-33} (guía de remisión) y \id{RF-34} (solicitud de insumos) son
mejoras deseables de bajo acoplamiento con el núcleo.

\textbf{\emph{Won't} en esta versión.} Se excluyen explícitamente y con argumento: la
\textbf{gestión de nómina y pago}, porque el sistema calcula el avance pero la liquidación
involucra información laboral y tributaria cuyo tratamiento excede el consentimiento obtenido
(\id{RD-04}, \id{RL-02}); la \textbf{facturación y contabilidad}, por no corresponder al manejo del
cultivo; la \textbf{comercialización de la producción}, que se negocia fuera del sistema; la
\textbf{integración directa con el sistema de la planta extractora}, porque pertenece a una
organización externa con sistema propio y su interfaz no está disponible (\id{EV-04}); y la
\textbf{operación multi-finca}, que el cliente administra pero que no forma parte del alcance
acordado para la v1 (§1.2.2).

\subsubsection{Modelo de Kano}

La clasificación se sustenta en la distancia entre la valoración declarada en el cuestionario y la
expectativa implícita del proceso actual.

\begin{longtable}{|L{2.8cm}|L{5.4cm}|L{7.0cm}|}
\hline
\rowcolor{grisclaro}\textbf{Categoría} & \textbf{Requisitos} & \textbf{Fundamento} \\
\hline
\endhead
Básica & \id{RF-01}, \id{RF-02}, \id{RF-03}, \id{RF-04} & Su ausencia genera rechazo, su presencia no genera entusiasmo. El registro de labor ya existe en papel: el sistema debe igualarlo, no impresiona por hacerlo. \\ \hline
De desempeño & \id{RF-19}, \id{RF-26}, \id{RF-28}, \id{RF-29}, \id{RF-13} & La satisfacción crece de forma proporcional a la calidad de ejecución. Valoración intermedia en el cuestionario ($3{,}95$ para el registro de labores). \\ \hline
\rowcolor{grisclaro}De deleite & \id{RF-07}, \id{RF-08}, \id{RF-12}, \id{RF-21} & Generan entusiasmo desproporcionado respecto de la expectativa. Son las funciones mejor valoradas ($4{,}27$ y $4{,}31$) y las únicas mencionadas de forma espontánea por los cinco perfiles entrevistados. \\ \hline
Indiferente & \id{RF-33}, \id{RF-34} & No inciden en la satisfacción percibida por la persona usuaria final. \\ \hline
De rechazo potencial & \id{RF-15} & El rastreo de recorridos puede generar insatisfacción: el 25{,}8\% del personal manifiesta reservas (\id{EV-12}). Se mitiga con consentimiento revocable (\id{RL-03}). \\ \hline
\caption{Clasificación de los requisitos según el modelo de Kano}
\end{longtable}

\noindent
La clasificación de \id{RF-15} como \emph{requisito de rechazo potencial} es deliberada y
constituye un hallazgo. La funcionalidad fue solicitada por la administración (\id{EV-01},
\id{EV-02}) y por la jefatura de polinización (\id{EV-03}), pero una de cada cuatro personas
destinatarias expresa reservas sobre ella. Kano permite representar esta tensión: un requisito puede
ser valioso para quien lo solicita y adverso para quien lo padece. La mitigación no es técnica sino
de diseño de la relación: consentimiento específico, revocable y sin consecuencia laboral.

\subsubsection{Valor de negocio mediante WSJF}

Se aplica el método del Trabajo Ponderado Más Corto Primero:
\[
\mathrm{WSJF} = \frac{\text{Valor de negocio} + \text{Criticidad temporal} + \text{Reducción de riesgo}}{\text{Tamaño del trabajo}}
\]
Los cuatro componentes se puntúan en escala Fibonacci (1, 2, 3, 5, 8, 13). La tabla presenta los
diez requisitos evaluados de mayor valor de negocio, ordenados de forma descendente por WSJF.

\begin{longtable}{|C{1.4cm}|C{1.3cm}|C{1.3cm}|C{1.3cm}|C{1.3cm}|C{1.4cm}|L{5.4cm}|}
\hline
\rowcolor{grisclaro}\textbf{RF} & \textbf{Valor} & \textbf{Crit. temp.} & \textbf{Red. riesgo} & \textbf{Tamaño} & \textbf{WSJF} & \textbf{Justificación} \\
\hline
\endhead
\id{RF-35} & 5 & 8 & 8 & 3 & \textbf{7,00} & Bajo costo y alto impacto: sin él, el 11,3\% del personal queda excluido del sistema (\id{EV-12}). \\ \hline
\id{RF-21} & 13 & 13 & 8 & 8 & \textbf{4,25} & Evita la penalización económica del lote; la temporada de baja producción (jul.--oct.) es cuando más fruta verde se corta (\id{EV-04}). \\ \hline
\id{RF-22} & 13 & 13 & 8 & 8 & \textbf{4,25} & Actúa antes del despacho, momento en que la pérdida aún es evitable. \\ \hline
\id{RF-04} & 8 & 8 & 5 & 5 & \textbf{4,20} & Base de todo el registro operativo; sin él no hay datos para ninguna otra función. \\ \hline
\id{RF-12} & 8 & 8 & 5 & 5 & \textbf{4,20} & Función mejor valorada del cuestionario (4,31/5). \\ \hline
\id{RF-01} & 5 & 8 & 8 & 5 & \textbf{4,20} & Precondición técnica y legal de todo lo demás (\id{RL-01}). \\ \hline
\id{RF-07} & 13 & 8 & 8 & 8 & \textbf{3,62} & Mayor dificultad declarada por el personal (40,3\%, \id{EV-12}). \\ \hline
\id{RF-28} & 8 & 5 & 5 & 5 & \textbf{3,60} & Habilita el cálculo de cumplimiento y de remuneración por avance. \\ \hline
\id{RF-14} & 13 & 8 & 5 & 8 & \textbf{3,25} & La polinización determina la producción; el conteo manual afecta el rendimiento (\id{EV-03}). \\ \hline
\id{RF-08} & 8 & 5 & 5 & 8 & \textbf{2,25} & Alto valor técnico, pero el diagnóstico nutricional admite mayor demora que el sanitario. \\ \hline

\caption{Valor de negocio mediante WSJF (diez requisitos de mayor prioridad)}
\end{longtable}

\noindent
\textbf{Lectura del resultado.} \id{RF-35} y \id{RF-10} encabezan el orden de implementación con un
WSJF de 7,00 cada uno,
pese a no ser el requisito de mayor valor de negocio. La razón es su tamaño reducido frente a un
impacto alto: se trata de una variante del formulario de registro que evita excluir a una de cada
nueve personas destinatarias. \id{RF-10} alcanza la misma posición por una razón distinta: es
prerrequisito de \id{RF-21}, dado que el umbral de desprendimiento se obtiene de la variedad del
lote. El método revela así una dependencia técnica que ni MoSCoW ni Kano hacen visible.

Ambos casos ilustran la utilidad del método frente a una priorización basada únicamente en valor
percibido, que los habría situado muy por debajo. La lectura inversa también es informativa:
\id{RF-07} y \id{RF-08}, las dos funciones mejor valoradas por las personas usuarias, obtienen WSJF
de 3,62 y 2,25 por su elevado tamaño de trabajo. Esto no altera su condición de \emph{Must}, pero
indica que el desarrollo debe iniciarse por la base operativa mientras los componentes de
inteligencia artificial se construyen en paralelo. El archivo
\id{04\_Trazabilidad/}\allowbreak\id{priorizacion\_moscow\_kano.csv} contiene la tabla completa para los 35
requisitos funcionales.

% ---------------------------------------------------------------------
\subsection{Matriz de trazabilidad extendida}

La matriz enlaza, para cada elemento, la cadena completa
Ley $\rightarrow$ Objetivo $\rightarrow$ Interesado $\rightarrow$ Evidencia $\rightarrow$
Requisito $\rightarrow$ Caso de uso $\rightarrow$ Historia de usuario $\rightarrow$
Criterio de aceptación $\rightarrow$ Componente $\rightarrow$ Mockup.

La versión completa, con las trece columnas exigidas por la §8.8 de la guía, reside en
\id{04\_Trazabilidad/}\allowbreak\id{matriz\_trazabilidad.csv} y contiene \textbf{52 filas}. La tabla siguiente
reproduce una vista reducida con las columnas de mayor densidad informativa; se conservan las 48
filas para permitir la verificación de cobertura. Las cuatro últimas corresponden a los requisitos derivados de la hoja de liquidación semanal (\id{EV-13}).

\begin{longtable}{|C{0.9cm}|C{1.5cm}|C{1.6cm}|C{1.5cm}|C{1.3cm}|C{1.3cm}|C{1.3cm}|C{1.5cm}|}
\hline
\rowcolor{grisclaro}\textbf{\#} & \textbf{Interesado} & \textbf{ID-EV} & \textbf{ID-RF} & \textbf{Tipo} & \textbf{ID-CU} & \textbf{ID-HU} & \textbf{ID-Mockup} \\
\hline
\endhead
1 & Administrador & EV-01, EV-09 & RF-01 & RF & CU-11 & HU-01 & MU-01 \\ \hline
2 & Administrador & EV-01, EV-02 & RF-02 & RF & CU-12 & HU-02 & MU-05 \\ \hline
3 & Administrador & EV-01, EV-11 & RF-03 & RF & CU-13 & HU-03 & --- \\ \hline
4 & Trabajador & EV-05, EV-06 & RF-04 & RF & CU-04 & HU-04 & MU-04 \\ \hline
5 & Trabajador & EV-11 & RF-04 & RF & CU-04 & HU-04 & MU-04 \\ \hline
6 & Técnico fitosan. & EV-01, EV-02 & RF-05 & RF & CU-15 & HU-05 & MU-05 \\ \hline
7 & Administrador & EV-01 & RF-06 & RF & CU-12 & --- & MU-05 \\ \hline
8 & Asistente admin. & EV-01, EV-02 & RF-07 & RF & CU-01 & HU-06 & MU-03 \\ \hline
9 & Asistente admin. & EV-07 & RF-07 & RF & CU-01 & HU-06 & MU-03 \\ \hline
10 & Asesor técnico & EV-02 & RF-08 & RF & CU-02 & HU-07 & MU-03 \\ \hline
11 & Asesor técnico & EV-04 & RF-08 & RF & CU-02 & HU-07 & MU-03 \\ \hline
12 & Técnico fitosan. & EV-01, EV-02 & RF-09 & RF & CU-01 & --- & MU-03 \\ \hline
13 & Asesor técnico & EV-04 & RF-10 & RF & CU-12 & HU-08 & MU-05 \\ \hline
14 & Asesor técnico & EV-02 & RF-11 & RF & CU-02 & --- & MU-05 \\ \hline
15 & Administrador & EV-01, EV-02 & RF-12 & RF & CU-16 & HU-09 & MU-08 \\ \hline
16 & Trabajador & EV-12 & RF-12 & RF & CU-16 & HU-09 & MU-08 \\ \hline
17 & Jefe polinización & EV-03 & RF-13 & RF & CU-03 & HU-10 & MU-04 \\ \hline
18 & Jefe polinización & EV-03 & RF-14 & RF & CU-03 & HU-11 & MU-06 \\ \hline
19 & Administrador & EV-01, EV-03 & RF-15 & RF & CU-03 & --- & MU-06 \\ \hline
20 & Administrador & EV-01, EV-02 & RF-16 & RF & CU-05 & --- & MU-07 \\ \hline
21 & Administrador & EV-01, EV-02 & RF-17 & RF & CU-05 & --- & MU-05 \\ \hline
22 & Administrador & EV-01, EV-02 & RF-18 & RF & CU-17 & HU-12 & MU-07 \\ \hline
23 & Administrador & EV-01, EV-11 & RF-19 & RF & CU-05 & HU-13 & MU-07 \\ \hline
24 & Administrador & EV-01 & RF-20 & RF & CU-05 & --- & MU-05 \\ \hline
25 & Téc. extractora & EV-04 & RF-21 & RF & CU-06 & HU-14 & MU-03 \\ \hline
26 & Téc. extractora & EV-08 & RF-21 & RF & CU-06 & HU-14 & MU-03 \\ \hline
27 & Administrador & EV-04, EV-08 & RF-22 & RF & CU-06 & HU-15 & MU-08 \\ \hline
28 & Administrador & EV-04 & RF-23 & RF & CU-10 & --- & MU-07 \\ \hline
29 & Téc. extractora & EV-04, EV-08 & RF-24 & RF & CU-10 & --- & MU-07 \\ \hline
30 & Téc. extractora & EV-04 & RF-25 & RF & CU-10 & --- & MU-07 \\ \hline
31 & Administrador & EV-11 & RF-26 & RF & CU-07 & HU-16 & MU-02 \\ \hline
32 & Administrador & EV-11 & RF-27 & RF & CU-07 & --- & MU-02 \\ \hline
33 & Trabajador & EV-05, EV-11 & RF-28 & RF & CU-04 & HU-17 & MU-04 \\ \hline
34 & Trabajador & EV-06, EV-07 & RF-29 & RF & CU-08 & --- & MU-02 \\ \hline
35 & Asistente admin. & EV-07 & RF-30 & RF & CU-09 & HU-18 & MU-06 \\ \hline
36 & Asistente admin. & EV-07 & RF-31 & RF & CU-09 & --- & MU-08 \\ \hline
37 & Téc. extractora & EV-04 & RF-32 & RF & CU-17 & --- & MU-07 \\ \hline
38 & Téc. extractora & EV-04 & RF-33 & RF & CU-10 & --- & MU-07 \\ \hline
39 & Administrador & EV-11 & RF-34 & RF & CU-07 & --- & MU-02 \\ \hline
40 & Jefe de campo & EV-12, EV-05 & RF-35 & RF & CU-04 & HU-19 & MU-04 \\ \hline
41 & Téc. extractora & EV-04 & RNF-01 & RNF & CU-06 & HU-14 & MU-03 \\ \hline
42 & Trabajador & EV-07, EV-12 & RNF-11 & RNF & CU-04 & HU-04 & MU-04 \\ \hline
43 & Jefe polinización & EV-03 & RNF-05 & RNF & CU-03 & HU-11 & MU-06 \\ \hline
44 & Trabajador & EV-12 & RNF-08 & RNF & CU-04 & HU-04 & MU-04 \\ \hline
45 & Trabajador & EV-12 & RNF-14 & RNF & CU-03 & HU-11 & MU-06 \\ \hline
46 & Asistente admin. & EV-01, EV-07 & RNF-16 & RNF & CU-18 & HU-06 & MU-03 \\ \hline
47 & Trabajador & EV-12 & RD-10 & RD & CU-04 & HU-19 & MU-04 \\ \hline
48 & Téc. extractora & EV-08 & RD-07 & RD & CU-06 & HU-14 & MU-03 \\ \hline
49 & Administrador & EV-13 & RF-36 & RF & CU-08 & --- & MU-02 \\ \hline
50 & Administrador & EV-13, EV-11 & RF-37 & RF & CU-07 & --- & MU-02 \\ \hline
51 & Administrador & EV-13 & RF-38 & RF & CU-07 & --- & MU-02 \\ \hline
52 & Administrador & EV-13 & RF-39 & RF & CU-07 & --- & MU-02 \\ \hline
\caption{Matriz de trazabilidad extendida (vista reducida, 52 filas)}
\end{longtable}

\subsubsection{Diagnóstico de cobertura}

\begin{longtable}{|L{5.6cm}|C{2.4cm}|L{6.8cm}|}
\hline
\rowcolor{grisclaro}\textbf{Indicador} & \textbf{Resultado} & \textbf{Interpretación} \\
\hline
\endhead
Requisitos funcionales trazados a evidencia & 39 / 39 (100\%) & Ningún requisito carece de fuente verificable. \\ \hline
Requisitos \emph{Must} con caso de uso & 19 / 19 (100\%) & Se corrige la deficiencia de la Entrega 2, donde \id{RF-17} y \id{RF-20} carecían de caso de uso. \\ \hline
Requisitos \emph{Must} con historia y criterio & 19 / 19 (100\%) & \\ \hline
Evidencias que originan al menos un requisito & 12 / 13 (92,3\%) & \id{EV-09} (observación de campo) aporta contexto y respalda requisitos junto con otras fuentes, sin originar ninguno de forma exclusiva. \\ \hline
Requisitos sustentados en una sola evidencia & 9 / 39 (23,1\%) & Véase el análisis de elementos huérfanos. \\ \hline
Requisitos con mockup asociado & 36 / 39 (92,3\%) & \id{RF-03}, \id{RF-27} y \id{RF-34} operan sobre pantallas no prototipadas en esta entrega. \\ \hline
\caption{Diagnóstico de cobertura de la matriz de trazabilidad}
\end{longtable}

\subsubsection{Elementos huérfanos y vacíos de cobertura}

Se declaran de forma explícita las debilidades detectadas, conforme al criterio de completitud de la
guía.

\begin{enumerate}[nosep]
  \item \textbf{Requisitos con fuente única.} \id{RF-11}, \id{RF-23}, \id{RF-25}, \id{RF-32},
        \id{RF-33} y \id{RF-34} se sustentan en una sola evidencia. Los cuatro derivados de
        \id{EV-04} conciernen a la relación con la planta extractora y no pudieron triangularse con
        una segunda fuente independiente de esa organización, dado que \id{EV-08} corresponde a la
        misma sesión de campo. Se declara como limitación.
  \item \textbf{Requisitos sin prototipo de interfaz.} \id{RF-03}, \id{RF-27} y \id{RF-34}. Los tres
        son de prioridad no \emph{Must}, por lo que la ausencia no incumple la exigencia de
        prototipar las pantallas \emph{Must}.
  \item \textbf{Ausencia de \emph{Won't} funcionales.} Las exclusiones se enuncian a nivel de
        producto y no como requisitos individuales descartados. Se argumentan en la §5.3.1.
  \item \textbf{Evidencia de contexto.} \id{EV-09} no origina requisitos de forma exclusiva. Se
        conserva en el catálogo por su valor de caracterización del proceso actual.
\end{enumerate}

% =====================================================================

\newpage
\section{Priorización y trazabilidad extendida}

\subsection{Clasificación de requisitos}

El documento especifica 75 elementos de requisito, distribuidos como sigue.

\begin{longtable}{|L{4.2cm}|C{2.0cm}|L{8.4cm}|}
\hline
\rowcolor{grisclaro}\textbf{Tipo} & \textbf{Cantidad} & \textbf{Identificadores} \\
\hline
\endhead
Requisito funcional (RF) & 39 & \id{RF-01} a \id{RF-39} \\ \hline
Requisito no funcional (RNF) & 18 & \id{RNF-01} a \id{RNF-18} \\ \hline
Restricción de diseño (RD) & 10 & \id{RD-01} a \id{RD-10} \\ \hline
Requisito legal (RL) & 8 & \id{RL-01} a \id{RL-08} \\ \hline
\rowcolor{grisclaro}\textbf{Total} & \textbf{75} & \\ \hline
\caption{Clasificación de los requisitos por tipo}
\end{longtable}

\noindent
La distribución MoSCoW de los requisitos funcionales es de 21 \emph{Must}, 16 \emph{Should} y 2
\emph{Could}. Los cuatro requisitos incorporados a partir de la hoja de liquidación semanal
(\id{RF-36} a \id{RF-39}) se clasifican como \emph{Must} los dos primeros ---el catálogo de tarifas
es prerrequisito del cálculo de remuneración y la liquidación es el cierre del ciclo semanal--- y
como \emph{Should} los dos restantes. No hay requisitos funcionales clasificados como \emph{Won't}: las exclusiones de
alcance se enuncian a nivel de producto en la §1.2.4 y se argumentan en la §5.3.4.

% ---------------------------------------------------------------------
\subsection{Historias de usuario de los requisitos \emph{Must}}

Conforme a la §3.3 de la guía de la Entrega 3, cada requisito funcional de prioridad \emph{Must}
cuenta con una historia de usuario en formato Connextra, con los criterios INVEST verificados y un
escenario de aceptación redactado en Gherkin. Se especifican 19 historias (\id{HU-01} a
\id{HU-19}), con sus criterios de aceptación (\id{CA-01} a \id{CA-19}).

\begin{longtable}{|C{1.3cm}|C{1.3cm}|L{12.4cm}|}
\hline
\rowcolor{grisclaro}\textbf{HU} & \textbf{RF} & \textbf{Historia de usuario (Connextra) y escenario de aceptación (Gherkin)} \\
\hline
\endhead

\id{HU-01} & \id{RF-01} & \textbf{Como} persona trabajadora del cultivo, \textbf{quiero} acceder al sistema con mis credenciales, \textbf{para} usar únicamente las funciones que corresponden a mi labor. \newline \emph{\id{CA-01}:} \textbf{Dado} que existe una cuenta con rol ``trabajador agrícola'', \textbf{cuando} inicia sesión, \textbf{entonces} el menú no presenta ninguna opción de administración \textbf{y} un intento de acceso directo a esa ruta es rechazado. \\ \hline

\id{HU-02} & \id{RF-02} & \textbf{Como} administrador general, \textbf{quiero} registrar los lotes con su variedad y superficie, \textbf{para} poder asociarles labores, diagnósticos y cosechas. \newline \emph{\id{CA-02}:} \textbf{Dado} un lote recién registrado con variedad asignada, \textbf{cuando} se abre el formulario de registro de labor, \textbf{entonces} el lote aparece como opción seleccionable. \\ \hline

\id{HU-03} & \id{RF-03} & \textbf{Como} administrador general, \textbf{quiero} organizar al personal por equipos funcionales, \textbf{para} asignar labores y evaluar el desempeño por cuadrilla. \newline \emph{\id{CA-03}:} \textbf{Dado} un trabajador asignado al equipo de polinización, \textbf{cuando} se consulta el equipo, \textbf{entonces} la persona figura en él \textbf{y} es asignable a labores de ese tipo. \\ \hline

\id{HU-04} & \id{RF-04} & \textbf{Como} trabajador agrícola, \textbf{quiero} registrar la labor que ejecuté en el lote, \textbf{para} que quede constancia sin depender de una libreta de papel. \newline \emph{\id{CA-04}:} \textbf{Dado} que ejecuté chapia en el lote 3, \textbf{cuando} registro la labor con su cantidad, \textbf{entonces} queda vinculada al lote, la fecha y mi identificación \textbf{y} aparece en el historial del lote. \\ \hline

\id{HU-05} & \id{RF-05} & \textbf{Como} técnico fitosanitario, \textbf{quiero} registrar el monitoreo sanitario del lote, \textbf{para} dejar trazabilidad de las revisiones y de las incidencias detectadas. \newline \emph{\id{CA-05}:} \textbf{Dado} un monitoreo con incidencia marcada, \textbf{cuando} se filtran los lotes con incidencia abierta, \textbf{entonces} ese lote aparece en el resultado. \\ \hline

\id{HU-06} & \id{RF-07} & \textbf{Como} asistente de administración, \textbf{quiero} fotografiar una palma afectada y obtener el diagnóstico, \textbf{para} identificar la plaga sin depender de conocimiento especializado. \newline \emph{\id{CA-06}:} \textbf{Dado} que capturo la imagen de una palma con afectación conocida, \textbf{cuando} solicito el análisis, \textbf{entonces} el sistema devuelve la clase, el nivel de confianza y una explicación de no más de 60 palabras \textbf{y} registra el diagnóstico con lote y fecha. \\ \hline

\id{HU-07} & \id{RF-08} & \textbf{Como} asesor técnico, \textbf{quiero} fotografiar una hoja y conocer la deficiencia nutricional, \textbf{para} decidir qué aplicar sin esperar el análisis de laboratorio. \newline \emph{\id{CA-07}:} \textbf{Dado} el lote con variedad asignada, \textbf{cuando} envío la imagen de una hoja con deficiencia conocida, \textbf{entonces} el sistema indica el elemento deficiente, propone una corrección \textbf{y} explica qué signo foliar la sustenta. \\ \hline

\id{HU-08} & \id{RF-10} & \textbf{Como} asesor técnico, \textbf{quiero} que el sistema conozca la variedad de cada lote, \textbf{para} que los umbrales de diagnóstico correspondan al material sembrado. \newline \emph{\id{CA-08}:} \textbf{Dado} un lote de variedad \emph{guineensis}, \textbf{cuando} se clasifica un racimo, \textbf{entonces} el umbral aplicado es de 3 frutos desprendidos; \textbf{y dado} un lote híbrido, \textbf{entonces} el umbral aplicado es de 5. \\ \hline

\id{HU-09} & \id{RF-12} & \textbf{Como} administrador general, \textbf{quiero} recibir alertas ante una afectación detectada, \textbf{para} intervenir antes de que se propague. \newline \emph{\id{CA-09}:} \textbf{Dado} un diagnóstico positivo registrado, \textbf{cuando} se completa el análisis, \textbf{entonces} se crea una alerta con lote, tipo, criticidad y estado abierto \textbf{y} resulta visible para el rol responsable. \\ \hline

\id{HU-10} & \id{RF-13} & \textbf{Como} jefe de polinización, \textbf{quiero} registrar el estado de la antesis y la verificación visual, \textbf{para} controlar la calidad de la labor que determina la producción. \newline \emph{\id{CA-10}:} \textbf{Dado} un lote en etapa de polinización, \textbf{cuando} se registra la jornada con estado de antesis, \textbf{entonces} el registro queda asociado al lote, la fecha y la persona \textbf{y} es consultable por fecha. \\ \hline

\id{HU-11} & \id{RF-14} & \textbf{Como} polinizador, \textbf{quiero} que el conteo de flores se calcule solo, \textbf{para} no perder rendimiento anotando planta por planta. \newline \emph{\id{CA-11}:} \textbf{Dado} que trabajé la jornada con rastreo activo, \textbf{cuando} cierro la jornada, \textbf{entonces} el sistema calcula el total de flores a partir de las marcaciones \textbf{y} no me solicita ningún registro manual por planta. \\ \hline

\id{HU-12} & \id{RF-18} & \textbf{Como} administrador general, \textbf{quiero} estimar la producción del período, \textbf{para} planificar la cosecha y la mano de obra necesaria. \newline \emph{\id{CA-12}:} \textbf{Dado} un censo de flores registrado y la variedad del lote, \textbf{cuando} solicito la estimación, \textbf{entonces} el sistema devuelve el tonelaje estimado del período \textbf{y} el supuesto de rendimiento empleado. \\ \hline

\id{HU-13} & \id{RF-19} & \textbf{Como} administrador general, \textbf{quiero} generar el consolidado semanal, \textbf{para} tomar decisiones sin recopilar reportes dispersos de mensajería. \newline \emph{\id{CA-13}:} \textbf{Dado} que existen registros en la semana, \textbf{cuando} genero el reporte, \textbf{entonces} obtengo labores, avance, producción y alertas del período \textbf{y} puedo exportarlo a PDF o a hoja de cálculo. \\ \hline

\id{HU-14} & \id{RF-21} & \textbf{Como} cosechador, \textbf{quiero} saber si un racimo está en punto de corte, \textbf{para} no cortar fruta verde que será penalizada. \newline \emph{\id{CA-14}:} \textbf{Dado} un lote híbrido, \textbf{cuando} fotografío la parte basal de un racimo con 4 frutos desprendidos, \textbf{entonces} el sistema lo clasifica como verde \textbf{y} explica que el umbral de la variedad es de 5. \\ \hline

\id{HU-15} & \id{RF-22} & \textbf{Como} administrador general, \textbf{quiero} conocer el porcentaje estimado de fruta verde antes de despachar, \textbf{para} evitar la penalización del lote en la planta extractora. \newline \emph{\id{CA-15}:} \textbf{Dado} un umbral de penalización del 3\%, \textbf{cuando} el porcentaje estimado del lote alcanza ese valor, \textbf{entonces} el sistema emite alerta preventiva \textbf{y} desaconseja el despacho. \\ \hline

\id{HU-16} & \id{RF-26} & \textbf{Como} administrador general, \textbf{quiero} registrar el presupuesto semanal de labores, \textbf{para} contrastar lo planificado con lo efectivamente ejecutado. \newline \emph{\id{CA-16}:} \textbf{Dado} un plan con meta de 1.123 coronas en el lote 3, \textbf{cuando} cierro la semana con 800 cumplidas, \textbf{entonces} el sistema reporta un cumplimiento del 71,2\% \textbf{y} un remanente de 323 coronas. \\ \hline

\id{HU-17} & \id{RF-28} & \textbf{Como} trabajador agrícola, \textbf{quiero} registrar mi avance en la unidad propia de cada labor, \textbf{para} que mi trabajo se contabilice como corresponde. \newline \emph{\id{CA-17}:} \textbf{Dado} que ejecuté 120 coronas, \textbf{cuando} registro el avance, \textbf{entonces} la cantidad se acumula a mi total semanal en la unidad ``coronas'' \textbf{y} el sistema rechaza una unidad incompatible con el tipo de labor. \\ \hline

\id{HU-18} & \id{RF-30} & \textbf{Como} asistente de administración, \textbf{quiero} reportar una afectación con foto y ubicación, \textbf{para} poder volver al punto exacto en mi siguiente visita. \newline \emph{\id{CA-18}:} \textbf{Dado} que reporto una incidencia con GPS activo, \textbf{cuando} la consulto días después, \textbf{entonces} recupero su fotografía y sus coordenadas \textbf{y} la incidencia figura en mi lista de pendientes de verificación. \\ \hline

\id{HU-19} & \id{RF-35} & \textbf{Como} jefe de campo, \textbf{quiero} registrar el avance de quien no tiene teléfono, \textbf{para} que ninguna persona quede sin registrar su trabajo. \newline \emph{\id{CA-19}:} \textbf{Dado} un trabajador sin dispositivo propio, \textbf{cuando} registro su avance de forma delegada, \textbf{entonces} el avance se atribuye a esa persona en su consolidado \textbf{y} la bitácora conserva mi identificación como registrador. \\ \hline
\caption{Historias de usuario y criterios de aceptación de los requisitos \emph{Must}}
\end{longtable}

\subsubsection{Verificación de los criterios INVEST}

Las diecinueve historias fueron evaluadas contra los seis criterios INVEST. La tabla siguiente
reporta el resultado y documenta las tres historias que requirieron ajuste.

\begin{longtable}{|L{2.4cm}|C{1.6cm}|L{11.0cm}|}
\hline
\rowcolor{grisclaro}\textbf{Criterio} & \textbf{Cumple} & \textbf{Observación} \\
\hline
\endhead
Independiente & 16 / 19 & \id{HU-11} depende de \id{HU-10}, \id{HU-15} de \id{HU-14} y \id{HU-17} de \id{HU-04}. Se documenta la dependencia en la matriz en lugar de forzar la independencia, dado que responde a la secuencia real del proceso. \\ \hline
Negociable & 19 / 19 & Ninguna historia prescribe solución técnica; los umbrales son parametrizables (\id{RD-08}). \\ \hline
Valiosa & 19 / 19 & Cada historia declara el beneficio esperado y se traza a evidencia de campo. \\ \hline
Estimable & 19 / 19 & Todas admiten estimación en la escala Fibonacci empleada en el cálculo WSJF de la §5.3.3. \\ \hline
Pequeña & 17 / 19 & \id{HU-06} y \id{HU-07} superan el tamaño recomendado por incluir el componente de IA; se sugiere descomponerlas en captura, inferencia y explicación durante la planificación de la implementación. \\ \hline
Verificable & 19 / 19 & Todas cuentan con escenario Gherkin con criterio observable. \\ \hline
\caption{Verificación de los criterios INVEST sobre las historias de usuario}
\end{longtable}

% ---------------------------------------------------------------------
\subsection{Priorización combinada}

La priorización emplea tres instrumentos complementarios: MoSCoW para el compromiso de alcance, el
modelo de Kano para la relación con la satisfacción del cliente, y WSJF para el orden de
implementación.

\subsubsection{Justificación de la clasificación MoSCoW}

\textbf{\emph{Must} (19 RF).} Corresponden a tres condiciones. Primero, la \textbf{base operativa
indispensable}: sin autenticación por rol (\id{RF-01}), estructura de lotes (\id{RF-02}), personal
(\id{RF-03}) y registro de labor y avance (\id{RF-04}, \id{RF-28}, \id{RF-35}) el sistema no
sustituye al proceso manual que pretende reemplazar. Segundo, el \textbf{diferencial de valor}: el
análisis por imagen (\id{RF-07}, \id{RF-08}, \id{RF-21}) y las alertas (\id{RF-12}) son las
funciones mejor valoradas por el personal en el cuestionario ($4{,}31$ y $4{,}27$ sobre 5,
\id{EV-12}) y las señaladas de forma independiente por la administración (\id{EV-01}) y la asesoría
técnica (\id{EV-02}). Tercero, la \textbf{consecuencia económica directa}: la polinización determina
la producción de la finca (\id{EV-03}), su conteo condiciona la evaluación del trabajo
(\id{RF-13}, \id{RF-14}), y el porcentaje de fruta verde determina la penalización del lote en la
planta extractora (\id{RF-22}, \id{EV-04}).

\textbf{\emph{Should} (14 RF).} Aportan valor sustantivo pero el sistema entrega su propuesta de
valor sin ellos en la primera versión: seguimiento de etapas, ficha de plaga, análisis de
suelo, recorridos, evaluación de desempeño, clima, historial, registro del ticket de la extractora,
trazabilidad de la penalización, ventana corte-entrega, arrastre semanal, consulta de remuneración,
lista de pendientes y exportación hacia la extractora.

\textbf{\emph{Could} (2 RF).} \id{RF-33} (guía de remisión) y \id{RF-34} (solicitud de insumos) son
mejoras deseables de bajo acoplamiento con el núcleo.

\textbf{\emph{Won't} en esta versión.} Se excluyen explícitamente y con argumento: la
\textbf{gestión de nómina y pago}, porque el sistema calcula el avance pero la liquidación
involucra información laboral y tributaria cuyo tratamiento excede el consentimiento obtenido
(\id{RD-04}, \id{RL-02}); la \textbf{facturación y contabilidad}, por no corresponder al manejo del
cultivo; la \textbf{comercialización de la producción}, que se negocia fuera del sistema; la
\textbf{integración directa con el sistema de la planta extractora}, porque pertenece a una
organización externa con sistema propio y su interfaz no está disponible (\id{EV-04}); y la
\textbf{operación multi-finca}, que el cliente administra pero que no forma parte del alcance
acordado para la v1 (§1.2.2).

\subsubsection{Modelo de Kano}

La clasificación se sustenta en la distancia entre la valoración declarada en el cuestionario y la
expectativa implícita del proceso actual.

\begin{longtable}{|L{2.8cm}|L{5.4cm}|L{7.0cm}|}
\hline
\rowcolor{grisclaro}\textbf{Categoría} & \textbf{Requisitos} & \textbf{Fundamento} \\
\hline
\endhead
Básica & \id{RF-01}, \id{RF-02}, \id{RF-03}, \id{RF-04} & Su ausencia genera rechazo, su presencia no genera entusiasmo. El registro de labor ya existe en papel: el sistema debe igualarlo, no impresiona por hacerlo. \\ \hline
De desempeño & \id{RF-19}, \id{RF-26}, \id{RF-28}, \id{RF-29}, \id{RF-13} & La satisfacción crece de forma proporcional a la calidad de ejecución. Valoración intermedia en el cuestionario ($3{,}95$ para el registro de labores). \\ \hline
\rowcolor{grisclaro}De deleite & \id{RF-07}, \id{RF-08}, \id{RF-12}, \id{RF-21} & Generan entusiasmo desproporcionado respecto de la expectativa. Son las funciones mejor valoradas ($4{,}27$ y $4{,}31$) y las únicas mencionadas de forma espontánea por los cinco perfiles entrevistados. \\ \hline
Indiferente & \id{RF-33}, \id{RF-34} & No inciden en la satisfacción percibida por la persona usuaria final. \\ \hline
De rechazo potencial & \id{RF-15} & El rastreo de recorridos puede generar insatisfacción: el 25{,}8\% del personal manifiesta reservas (\id{EV-12}). Se mitiga con consentimiento revocable (\id{RL-03}). \\ \hline
\caption{Clasificación de los requisitos según el modelo de Kano}
\end{longtable}

\noindent
La clasificación de \id{RF-15} como \emph{requisito de rechazo potencial} es deliberada y
constituye un hallazgo. La funcionalidad fue solicitada por la administración (\id{EV-01},
\id{EV-02}) y por la jefatura de polinización (\id{EV-03}), pero una de cada cuatro personas
destinatarias expresa reservas sobre ella. Kano permite representar esta tensión: un requisito puede
ser valioso para quien lo solicita y adverso para quien lo padece. La mitigación no es técnica sino
de diseño de la relación: consentimiento específico, revocable y sin consecuencia laboral.

\subsubsection{Valor de negocio mediante WSJF}

Se aplica el método del Trabajo Ponderado Más Corto Primero:
\[
\mathrm{WSJF} = \frac{\text{Valor de negocio} + \text{Criticidad temporal} + \text{Reducción de riesgo}}{\text{Tamaño del trabajo}}
\]
Los cuatro componentes se puntúan en escala Fibonacci (1, 2, 3, 5, 8, 13). La tabla presenta los
diez requisitos evaluados de mayor valor de negocio, ordenados de forma descendente por WSJF.

\begin{longtable}{|C{1.4cm}|C{1.3cm}|C{1.3cm}|C{1.3cm}|C{1.3cm}|C{1.4cm}|L{5.4cm}|}
\hline
\rowcolor{grisclaro}\textbf{RF} & \textbf{Valor} & \textbf{Crit. temp.} & \textbf{Red. riesgo} & \textbf{Tamaño} & \textbf{WSJF} & \textbf{Justificación} \\
\hline
\endhead
\id{RF-35} & 5 & 8 & 8 & 3 & \textbf{7,00} & Bajo costo y alto impacto: sin él, el 11,3\% del personal queda excluido del sistema (\id{EV-12}). \\ \hline
\id{RF-21} & 13 & 13 & 8 & 8 & \textbf{4,25} & Evita la penalización económica del lote; la temporada de baja producción (jul.--oct.) es cuando más fruta verde se corta (\id{EV-04}). \\ \hline
\id{RF-22} & 13 & 13 & 8 & 8 & \textbf{4,25} & Actúa antes del despacho, momento en que la pérdida aún es evitable. \\ \hline
\id{RF-04} & 8 & 8 & 5 & 5 & \textbf{4,20} & Base de todo el registro operativo; sin él no hay datos para ninguna otra función. \\ \hline
\id{RF-12} & 8 & 8 & 5 & 5 & \textbf{4,20} & Función mejor valorada del cuestionario (4,31/5). \\ \hline
\id{RF-01} & 5 & 8 & 8 & 5 & \textbf{4,20} & Precondición técnica y legal de todo lo demás (\id{RL-01}). \\ \hline
\id{RF-07} & 13 & 8 & 8 & 8 & \textbf{3,62} & Mayor dificultad declarada por el personal (40,3\%, \id{EV-12}). \\ \hline
\id{RF-28} & 8 & 5 & 5 & 5 & \textbf{3,60} & Habilita el cálculo de cumplimiento y de remuneración por avance. \\ \hline
\id{RF-14} & 13 & 8 & 5 & 8 & \textbf{3,25} & La polinización determina la producción; el conteo manual afecta el rendimiento (\id{EV-03}). \\ \hline
\id{RF-08} & 8 & 5 & 5 & 8 & \textbf{2,25} & Alto valor técnico, pero el diagnóstico nutricional admite mayor demora que el sanitario. \\ \hline

\caption{Valor de negocio mediante WSJF (diez requisitos de mayor prioridad)}
\end{longtable}

\noindent
\textbf{Lectura del resultado.} \id{RF-35} y \id{RF-10} encabezan el orden de implementación con un
WSJF de 7,00 cada uno,
pese a no ser el requisito de mayor valor de negocio. La razón es su tamaño reducido frente a un
impacto alto: se trata de una variante del formulario de registro que evita excluir a una de cada
nueve personas destinatarias. \id{RF-10} alcanza la misma posición por una razón distinta: es
prerrequisito de \id{RF-21}, dado que el umbral de desprendimiento se obtiene de la variedad del
lote. El método revela así una dependencia técnica que ni MoSCoW ni Kano hacen visible.

Ambos casos ilustran la utilidad del método frente a una priorización basada únicamente en valor
percibido, que los habría situado muy por debajo. La lectura inversa también es informativa:
\id{RF-07} y \id{RF-08}, las dos funciones mejor valoradas por las personas usuarias, obtienen WSJF
de 3,62 y 2,25 por su elevado tamaño de trabajo. Esto no altera su condición de \emph{Must}, pero
indica que el desarrollo debe iniciarse por la base operativa mientras los componentes de
inteligencia artificial se construyen en paralelo. El archivo
\id{04\_Trazabilidad/}\allowbreak\id{priorizacion\_moscow\_kano.csv} contiene la tabla completa para los 35
requisitos funcionales.

% ---------------------------------------------------------------------
\subsection{Matriz de trazabilidad extendida}

La matriz enlaza, para cada elemento, la cadena completa
Ley $\rightarrow$ Objetivo $\rightarrow$ Interesado $\rightarrow$ Evidencia $\rightarrow$
Requisito $\rightarrow$ Caso de uso $\rightarrow$ Historia de usuario $\rightarrow$
Criterio de aceptación $\rightarrow$ Componente $\rightarrow$ Mockup.

La versión completa, con las trece columnas exigidas por la §8.8 de la guía, reside en
\id{04\_Trazabilidad/}\allowbreak\id{matriz\_trazabilidad.csv} y contiene \textbf{52 filas}. La tabla siguiente
reproduce una vista reducida con las columnas de mayor densidad informativa; se conservan las 48
filas para permitir la verificación de cobertura. Las cuatro últimas corresponden a los requisitos derivados de la hoja de liquidación semanal (\id{EV-13}).

\begin{longtable}{|C{0.9cm}|C{1.5cm}|C{1.6cm}|C{1.5cm}|C{1.3cm}|C{1.3cm}|C{1.3cm}|C{1.5cm}|}
\hline
\rowcolor{grisclaro}\textbf{\#} & \textbf{Interesado} & \textbf{ID-EV} & \textbf{ID-RF} & \textbf{Tipo} & \textbf{ID-CU} & \textbf{ID-HU} & \textbf{ID-Mockup} \\
\hline
\endhead
1 & Administrador & EV-01, EV-09 & RF-01 & RF & CU-11 & HU-01 & MU-01 \\ \hline
2 & Administrador & EV-01, EV-02 & RF-02 & RF & CU-12 & HU-02 & MU-05 \\ \hline
3 & Administrador & EV-01, EV-11 & RF-03 & RF & CU-13 & HU-03 & --- \\ \hline
4 & Trabajador & EV-05, EV-06 & RF-04 & RF & CU-04 & HU-04 & MU-04 \\ \hline
5 & Trabajador & EV-11 & RF-04 & RF & CU-04 & HU-04 & MU-04 \\ \hline
6 & Técnico fitosan. & EV-01, EV-02 & RF-05 & RF & CU-15 & HU-05 & MU-05 \\ \hline
7 & Administrador & EV-01 & RF-06 & RF & CU-12 & --- & MU-05 \\ \hline
8 & Asistente admin. & EV-01, EV-02 & RF-07 & RF & CU-01 & HU-06 & MU-03 \\ \hline
9 & Asistente admin. & EV-07 & RF-07 & RF & CU-01 & HU-06 & MU-03 \\ \hline
10 & Asesor técnico & EV-02 & RF-08 & RF & CU-02 & HU-07 & MU-03 \\ \hline
11 & Asesor técnico & EV-04 & RF-08 & RF & CU-02 & HU-07 & MU-03 \\ \hline
12 & Técnico fitosan. & EV-01, EV-02 & RF-09 & RF & CU-01 & --- & MU-03 \\ \hline
13 & Asesor técnico & EV-04 & RF-10 & RF & CU-12 & HU-08 & MU-05 \\ \hline
14 & Asesor técnico & EV-02 & RF-11 & RF & CU-02 & --- & MU-05 \\ \hline
15 & Administrador & EV-01, EV-02 & RF-12 & RF & CU-16 & HU-09 & MU-08 \\ \hline
16 & Trabajador & EV-12 & RF-12 & RF & CU-16 & HU-09 & MU-08 \\ \hline
17 & Jefe polinización & EV-03 & RF-13 & RF & CU-03 & HU-10 & MU-04 \\ \hline
18 & Jefe polinización & EV-03 & RF-14 & RF & CU-03 & HU-11 & MU-06 \\ \hline
19 & Administrador & EV-01, EV-03 & RF-15 & RF & CU-03 & --- & MU-06 \\ \hline
20 & Administrador & EV-01, EV-02 & RF-16 & RF & CU-05 & --- & MU-07 \\ \hline
21 & Administrador & EV-01, EV-02 & RF-17 & RF & CU-05 & --- & MU-05 \\ \hline
22 & Administrador & EV-01, EV-02 & RF-18 & RF & CU-17 & HU-12 & MU-07 \\ \hline
23 & Administrador & EV-01, EV-11 & RF-19 & RF & CU-05 & HU-13 & MU-07 \\ \hline
24 & Administrador & EV-01 & RF-20 & RF & CU-05 & --- & MU-05 \\ \hline
25 & Téc. extractora & EV-04 & RF-21 & RF & CU-06 & HU-14 & MU-03 \\ \hline
26 & Téc. extractora & EV-08 & RF-21 & RF & CU-06 & HU-14 & MU-03 \\ \hline
27 & Administrador & EV-04, EV-08 & RF-22 & RF & CU-06 & HU-15 & MU-08 \\ \hline
28 & Administrador & EV-04 & RF-23 & RF & CU-10 & --- & MU-07 \\ \hline
29 & Téc. extractora & EV-04, EV-08 & RF-24 & RF & CU-10 & --- & MU-07 \\ \hline
30 & Téc. extractora & EV-04 & RF-25 & RF & CU-10 & --- & MU-07 \\ \hline
31 & Administrador & EV-11 & RF-26 & RF & CU-07 & HU-16 & MU-02 \\ \hline
32 & Administrador & EV-11 & RF-27 & RF & CU-07 & --- & MU-02 \\ \hline
33 & Trabajador & EV-05, EV-11 & RF-28 & RF & CU-04 & HU-17 & MU-04 \\ \hline
34 & Trabajador & EV-06, EV-07 & RF-29 & RF & CU-08 & --- & MU-02 \\ \hline
35 & Asistente admin. & EV-07 & RF-30 & RF & CU-09 & HU-18 & MU-06 \\ \hline
36 & Asistente admin. & EV-07 & RF-31 & RF & CU-09 & --- & MU-08 \\ \hline
37 & Téc. extractora & EV-04 & RF-32 & RF & CU-17 & --- & MU-07 \\ \hline
38 & Téc. extractora & EV-04 & RF-33 & RF & CU-10 & --- & MU-07 \\ \hline
39 & Administrador & EV-11 & RF-34 & RF & CU-07 & --- & MU-02 \\ \hline
40 & Jefe de campo & EV-12, EV-05 & RF-35 & RF & CU-04 & HU-19 & MU-04 \\ \hline
41 & Téc. extractora & EV-04 & RNF-01 & RNF & CU-06 & HU-14 & MU-03 \\ \hline
42 & Trabajador & EV-07, EV-12 & RNF-11 & RNF & CU-04 & HU-04 & MU-04 \\ \hline
43 & Jefe polinización & EV-03 & RNF-05 & RNF & CU-03 & HU-11 & MU-06 \\ \hline
44 & Trabajador & EV-12 & RNF-08 & RNF & CU-04 & HU-04 & MU-04 \\ \hline
45 & Trabajador & EV-12 & RNF-14 & RNF & CU-03 & HU-11 & MU-06 \\ \hline
46 & Asistente admin. & EV-01, EV-07 & RNF-16 & RNF & CU-18 & HU-06 & MU-03 \\ \hline
47 & Trabajador & EV-12 & RD-10 & RD & CU-04 & HU-19 & MU-04 \\ \hline
48 & Téc. extractora & EV-08 & RD-07 & RD & CU-06 & HU-14 & MU-03 \\ \hline
49 & Administrador & EV-13 & RF-36 & RF & CU-08 & --- & MU-02 \\ \hline
50 & Administrador & EV-13, EV-11 & RF-37 & RF & CU-07 & --- & MU-02 \\ \hline
51 & Administrador & EV-13 & RF-38 & RF & CU-07 & --- & MU-02 \\ \hline
52 & Administrador & EV-13 & RF-39 & RF & CU-07 & --- & MU-02 \\ \hline
\caption{Matriz de trazabilidad extendida (vista reducida, 52 filas)}
\end{longtable}

\subsubsection{Diagnóstico de cobertura}

\begin{longtable}{|L{5.6cm}|C{2.4cm}|L{6.8cm}|}
\hline
\rowcolor{grisclaro}\textbf{Indicador} & \textbf{Resultado} & \textbf{Interpretación} \\
\hline
\endhead
Requisitos funcionales trazados a evidencia & 39 / 39 (100\%) & Ningún requisito carece de fuente verificable. \\ \hline
Requisitos \emph{Must} con caso de uso & 19 / 19 (100\%) & Se corrige la deficiencia de la Entrega 2, donde \id{RF-17} y \id{RF-20} carecían de caso de uso. \\ \hline
Requisitos \emph{Must} con historia y criterio & 19 / 19 (100\%) & \\ \hline
Evidencias que originan al menos un requisito & 12 / 13 (92,3\%) & \id{EV-09} (observación de campo) aporta contexto y respalda requisitos junto con otras fuentes, sin originar ninguno de forma exclusiva. \\ \hline
Requisitos sustentados en una sola evidencia & 9 / 39 (23,1\%) & Véase el análisis de elementos huérfanos. \\ \hline
Requisitos con mockup asociado & 36 / 39 (92,3\%) & \id{RF-03}, \id{RF-27} y \id{RF-34} operan sobre pantallas no prototipadas en esta entrega. \\ \hline
\caption{Diagnóstico de cobertura de la matriz de trazabilidad}
\end{longtable}

\subsubsection{Elementos huérfanos y vacíos de cobertura}

Se declaran de forma explícita las debilidades detectadas, conforme al criterio de completitud de la
guía.

\begin{enumerate}[nosep]
  \item \textbf{Requisitos con fuente única.} \id{RF-11}, \id{RF-23}, \id{RF-25}, \id{RF-32},
        \id{RF-33} y \id{RF-34} se sustentan en una sola evidencia. Los cuatro derivados de
        \id{EV-04} conciernen a la relación con la planta extractora y no pudieron triangularse con
        una segunda fuente independiente de esa organización, dado que \id{EV-08} corresponde a la
        misma sesión de campo. Se declara como limitación.
  \item \textbf{Requisitos sin prototipo de interfaz.} \id{RF-03}, \id{RF-27} y \id{RF-34}. Los tres
        son de prioridad no \emph{Must}, por lo que la ausencia no incumple la exigencia de
        prototipar las pantallas \emph{Must}.
  \item \textbf{Ausencia de \emph{Won't} funcionales.} Las exclusiones se enuncian a nivel de
        producto y no como requisitos individuales descartados. Se argumentan en la §5.3.1.
  \item \textbf{Evidencia de contexto.} \id{EV-09} no origina requisitos de forma exclusiva. Se
        conserva en el catálogo por su valor de caracterización del proceso actual.
\end{enumerate}
